\documentclass[11pt]{article}
\usepackage{tcee}
\addbibresource{3A03.bib}
\tcetema[Los economistas neoclásicos]{3.A.3}{Los economistas neoclásicos.}
\tcefecha{04/11/2026}
\begin{document}
\tceetitulo

\section*{Introducción}
\minutos{2,5}

\begin{itemize}
\item \textbf{\emph{Enganche}}:
  \begin{itemize}
  \item «Nadie sabe lo que vale el agua hasta que el pozo se seca», dice el refrán.\footnote{Es un refrán antiguo, también en inglés: \autor{Benjamin Franklin} lo recoge en su almanaque de 1746 («When the Well's dry, we know the Worth of Water») \parencite{Franklin1746PoorRichard}.}
  \item \autor{Adam Smith} se preguntaba por qué el agua, lo más útil que hay, apenas se puede cambiar por nada, mientras que por un diamante, que casi no tiene uso, se obtiene a menudo una cantidad enorme de otros bienes \parencite[libro~I, cap.~IV]{Smith1776Wealth}.
  \item Con un valor que dependía del coste de producción, la paradoja del agua y los diamantes no tenía buena respuesta.
  \item Entre 1871 y 1874, 3 economistas que no se conocían, en Manchester, Viena y Lausana, dieron la misma: lo que fija el valor de un bien no es su utilidad total, sino la de su \emph{última} unidad, la \concepto{utilidad marginal}.
  \item Es lo que dice el refrán: mientras el pozo está lleno, un vaso más o menos apenas importa; cuando se seca, el último vaso vale una fortuna. El agua es la misma; lo que cambia es la escasez, y con ella el valor de la última unidad.
  \item Detrás está la pregunta de esta exposición: \emph{¿qué determina el valor de las cosas: lo que cuesta producirlas o lo que las valoramos?}
  \end{itemize}

\item \textbf{\emph{Relevancia}}:
  \begin{itemize}
  \item \emph{Para la teoría}: la microeconomía que se estudia hoy está hecha con el análisis marginal de los neoclásicos: la demanda \vertema{3.A.8}, la producción y los costes \vertema{3.A.11}, los mercados \vertema{3.A.16}, el equilibrio general \vertema{3.A.21} y la economía del bienestar \vertema{3.A.22}.
  \item \emph{Para la política económica}: los precios como señales de escasez, el excedente del consumidor en el análisis coste-beneficio, el impuesto pigouviano sobre las externalidades, que es lo que son los impuestos sobre el carbono \vertema{3.A.23}, y el tipo de interés natural de \autor{Wicksell}, que los bancos centrales siguen estimando para orientar su política \parencite{HolstonLaubachWilliams2017Measuring} \vertema{3.A.7} \vertema{3.A.36}.
  \item \emph{Para el temario}: es la teoría que \autor{Keynes} llamará «clásica» y contra la que escribe la \emph{Teoría general} \vertema{3.A.4}.
  \item Además, como toda historia del pensamiento económico, enseña a no quedarse con el último manual y a ver por qué unas ideas llegan en unas épocas y no en otras: son 2 de los 4 motivos de \textcite{Schumpeter1954History}.\footnote{Los otros 2: la historia es fuente de ideas (principios que no fructificaron pueden orientar la investigación de hoy) y muestra cómo se desplazan las fronteras de la ciencia (cada época trae sus problemas y su forma de afrontarlos). La evolución del pensamiento económico resulta de la interacción de varios factores: el contexto económico de cada momento, el pensamiento anterior y los avances de otras disciplinas, como la filosofía, las matemáticas, la física o la biología. Conocerla permite entender las raíces del pensamiento actual, analizar los fenómenos económicos pasados y presentes y valorar los distintos programas de investigación.}
  \end{itemize}

\item \textbf{\emph{Contextualización}}:
  \begin{itemize}
  \item La historia del pensamiento económico avanza por \emph{revoluciones}: un \emph{paradigma} domina hasta que acumula anomalías que no sabe explicar y otro lo sustituye \parencite{Kuhn1962Structure} \vertema{3.A.1}. Si la marginalista fue una revolución en ese sentido es discutible: la nueva teoría tardó unos 20 años en imponerse, más un proceso que un vuelco \parencite{Blaug1972Marginal}.
  \item Los \emph{clásicos} \vertema{3.A.2}, de \autor{Smith} (1776) a \autor{John Stuart Mill} (1848), explicaban el valor por el coste de producción, en último término por el trabajo, y estudiaban el crecimiento y el reparto de la renta entre las clases sociales (salarios, beneficios y rentas de la tierra). \autor{Mill} daba la cuestión por cerrada: «felizmente, no queda nada en las leyes del valor que un escritor actual o futuro tenga que aclarar; la teoría está completa» \parencite[libro~III, cap.~I, §~1; traducción propia]{Mill1848Principles}.
  \item Quedaban, sin embargo, cabos sueltos:
    \begin{itemize}
    \item la paradoja del agua y los diamantes;
    \item los bienes que no se pueden reproducir (una estatua, un vino raro), cuyo valor depende solo de su escasez, como admitía el propio \textcite[cap.~I, secc.~I]{Ricardo1817Principles};
    \item y una teoría del valor trabajo de la que \autor{Marx} había sacado la explotación \vertema{3.A.2}.
    \end{itemize}
  \item El contexto económico también había cambiado. Es la segunda revolución industrial: entre 1870 y 1913 el PIB por habitante crece un 0,8~\% al año en el Reino Unido, frente al 1,6~\% de Alemania y el 1,75~\% de Estados Unidos, que lo supera desde 1880 \parencite{GGDC2024Maddison}. Lejos del estado estacionario que temían los clásicos, la pregunta urgente ya no es cómo crecer, sino cómo se asignan los recursos y se forman los precios.
  \item Y la economía se hace profesión: en 1885 \autor{Alfred Marshall} gana la cátedra de Cambridge y en 1890 impulsa la British Economic Association, que publica el \emph{Economic Journal} desde 1891 \parencite{HET2026Marshall,HET2026RES}.
  \item En 1871, \autor{William Stanley Jevons} publica en Londres \emph{The Theory of Political Economy} \parencite{Jevons1871Theory}; ese mismo año, \autor{Carl Menger} publica en Viena sus \emph{Grundsätze der Volkswirthschaftslehre} \parencite{Menger1871Grundsatze}, y en 1874 \autor{Léon Walras}, francés que enseñaba en Lausana, la primera entrega de sus \emph{Éléments d'économie politique pure} (la segunda es de 1877) \parencite{Walras1874Elements}. \autor{Marshall} sostuvo después que tenía estas ideas a finales de los años 60, aunque no publica sus \emph{Principles} hasta 1890 \parencite{Marshall1890Principles,HET2026Marshall}.
  \item La revolución no se hace de una vez. \textcite{Fonseca2026Neoclassicals} distingue 3 etapas, que ordenan esta exposición:
    \begin{enumerate}
    \item la \concepto{revolución marginalista} (1871-1874): \autor{Jevons}, \autor{Menger} y \autor{Walras} basan la teoría del intercambio en la utilidad marginal decreciente; con \autor{Marshall}, que la consolida en 1890, es el bloque~\ref{sec:revolucion};
    \item el \concepto{establecimiento de la teoría de la distribución según la productividad marginal} (1890-1894): \autor{John Bates Clark}, \autor{Philip Wicksteed} y \autor{Knut Wicksell} llevan el margen a la oferta y al reparto de la renta (\ref{sec:productividad});
    \item el \concepto{resurgimiento paretiano} (1934-1947): tras unas décadas en que la escuela pierde impulso, \autor{John Hicks}, \autor{Harold Hotelling}, \autor{Oskar Lange}, \autor{Maurice Allais} y \autor{Paul Samuelson} la rehacen con la utilidad ordinal de \autor{Pareto} y la nueva economía del bienestar (\ref{sec:paretiano}).
    \end{enumerate}
  \item Las 2 últimas etapas forman el bloque~\ref{sec:segunda}; entre una y otra, \autor{Fisher}, \autor{Wicksell}, \autor{Edgeworth}, \autor{Pareto} y \autor{Pigou} llevan el instrumental al interés, al dinero y al bienestar (\ref{sec:interes}).
  \item El análisis marginal lleva el centro de la disciplina a la microeconomía: a cómo asignan los mercados competitivos unos recursos escasos entre usos alternativos. Lo agregado no desaparece, pero queda en la teoría monetaria, separada de la del valor, y en la del ciclo.
  \item Para \textcite{ScrepantiZamagni2005Outline}, la revolución marginalista rompe con el método, los conceptos y las teorías de los clásicos.
  \item El nombre de \concepto{neoclásicos} lo pone \autor{Thorstein Veblen} en 1900, para subrayar la continuidad de \autor{Marshall} con los clásicos \parencite{Veblen1900Preconceptions,Aspromourgos1986Neoclassical}.\footnote{\autor{Veblen} es el fundador del institucionalismo norteamericano, una escuela muy crítica con la economía de los marginalistas, a la que reprochaba su visión estática y su individuo calculador.} Después, \autor{John Hicks} y \autor{George Stigler} \parencite*{Stigler1941Production} lo extienden a todos los marginalistas, que es el sentido que usamos aquí; \autor{Keynes} los llamará a todos «clásicos» \vertema{3.A.4}.
  \end{itemize}

\item \textbf{\emph{Problemática}}:
  \begin{itemize}
  \item ¿Qué cambia la revolución marginalista frente a los clásicos: la teoría del valor, el objeto o el método? ¿Fue de verdad una revolución?
  \item ¿Qué tienen en común las escuelas de Cambridge, Lausana y Viena, y qué las separa?
  \item ¿Cómo pasa el análisis marginal de la demanda a la producción y a la distribución de la renta?
  \item ¿Qué dijeron sobre el interés, el dinero y el bienestar?
  \item ¿Por qué la escuela pierde impulso a principios del siglo~XX, cómo se rehace en los años 30 y qué queda hoy de su edificio?
  \end{itemize}

\item \textbf{\emph{Estructura}}:
  \begin{itemize}
  \item La exposición sigue el orden en que se construyó la teoría, las 3 etapas de \autor{Fonseca}, del valor y la demanda a la oferta, la distribución y el bienestar:
  \begin{itemize}
  \item \concepto{I. Caracterización general de la escuela neoclásica} (\ref{sec:caracterizacion}): quiénes forman la escuela, de dónde viene, cuándo se desarrolla y con qué método, y en qué se distinguen sus 3 escuelas;
  \item \concepto{II. Revolución marginalista (1871-1874)} (\ref{sec:revolucion}): la 1.ª etapa, con los 4 fundadores, \autor{Jevons}, \autor{Menger}, \autor{Walras} y \autor{Marshall}, de la utilidad marginal a las tijeras de la oferta y la demanda;
  \item \concepto{III. Productividad marginal y resurgimiento paretiano} (\ref{sec:segunda}): la 2.ª etapa, en la que el margen llega a la producción y al reparto de la renta (\autor{Clark}, \autor{Wicksteed}, \autor{Wicksell}); el interés, el dinero y el bienestar (\autor{Fisher}, \autor{Edgeworth}, \autor{Pareto}, \autor{Pigou}), y la 3.ª etapa, en la que \autor{Hicks}, \autor{Lange}, \autor{Allais} y \autor{Samuelson} rehacen la teoría con la utilidad ordinal y los teoremas del bienestar.
  \end{itemize}
  \end{itemize}
\end{itemize}

\begin{esquema}[unbreakable]
\textbf{3.A.3. Los economistas neoclásicos}\\[2pt]
Introducción\esqtiempo{2,5}\\
I. Caracterización general de la escuela neoclásica\esqtiempo{5}\\
\hspace*{1.5em}1. Caracterización general de la escuela neoclásica según \autor{Schumpeter}\\
\hspace*{1.5em}2. Autores y rango histórico\\
\hspace*{1.5em}3. Método\\
II. Revolución marginalista (1871-1874)\esqtiempo{12,5}\\
\hspace*{1.5em}1. \autor{William Stanley Jevons} (1835-1882)\\
\hspace*{1.5em}2. \autor{Carl Menger} (1840-1921) y la escuela de Viena\\
\hspace*{1.5em}3. \autor{Léon Walras} (1834-1910) y la escuela de Lausana\\
\hspace*{1.5em}4. \autor{Alfred Marshall} (1842-1924) y la escuela de Cambridge\\
III. Productividad marginal y resurgimiento paretiano\esqtiempo{9}\\
\hspace*{1.5em}1. Establecimiento de la teoría de la distribución de la productividad marginal (1890-1899)\\
\hspace*{1.5em}2. Interés, dinero y bienestar\\
\hspace*{1.5em}3. Resurgimiento paretiano (1934-1947)\\
Conclusión\esqtiempo{1}
\end{esquema}

\begin{notaopositor}
Reparto del cante: introducción, \tiempo{2,5}; I, \tiempo{5} (I.1, \tiempo{1,5}; I.2, \tiempo{1,5}; I.3, \tiempo{2}); II, \tiempo{12,5} (II.1, \tiempo{2}; II.2, \tiempo{3,5}; II.3, \tiempo{2,5}; II.4, \tiempo{4,5}); III, \tiempo{9} (III.1, \tiempo{3,5}; III.2, \tiempo{3,5}; III.3, \tiempo{2}), y conclusión, \tiempo{1}. Las 3 etapas de \autor{Fonseca} se dicen en la contextualización y ordenan los bloques~\ref{sec:revolucion} y~\ref{sec:segunda}: conviene apuntarlas en la pizarra, pequeñas, como una línea del tiempo (1871-1874, 1890-1894, 1934-1947). Las cajas grises no se cantan. \autor{Marshall} no baja de \tiempo{4,5}: esta exposición es la de referencia de su equilibrio parcial y el tribunal espera ver el mercado de pescado en la pizarra. Las 3 escuelas de Cambridge, Lausana y Viena se presentan en el método (\ref{sec:metodo}) y vuelven con cada fundador. La pizarra se organiza en 3 zonas. A la izquierda, el refrán del enganche («nadie sabe lo que vale el agua hasta que el pozo se seca») y las pegatinas, cada una junto a su concepto: «lo poco agrada y lo mucho enfada» (la utilidad marginal decreciente), «cada loco con su tema» (la disputa sobre el método), «una mano lava la otra, y las dos, la cara» (las tijeras de \autor{Marshall}), «más vale un toma que dos te daré» (la impaciencia y el interés) y «a cada cual lo suyo» (la productividad marginal de \autor{Clark}). En el centro, los gráficos, que se borran al cambiar de bloque: el agua y los diamantes, la tabla de \autor{Menger} y el mercado de pescado. A la derecha, el esquema de las 3 escuelas, que se dibuja en el método y se llena con cada autor, y, recuadradas y conservadas desde que salen hasta el final, las 2 funciones clave: la segunda ley de \autor{Gossen} y el agotamiento del producto. Lo que no se canta (los anexos, las notas y los párrafos de profundización) sirve para las preguntas del tribunal.
\end{notaopositor}

% ==========================================================================
\section{Caracterización general de la escuela neoclásica}\label{sec:caracterizacion}
\minutos{5}

\begin{itemize}
\item Empezamos por la escuela en su conjunto: quiénes la forman, de dónde viene, cuándo se desarrolla y con qué método trabaja.
\item \emph{Conclusión anticipada}: la escuela neoclásica no tiene un líder, sino 3 escuelas nacionales, Cambridge, Lausana y Viena, que comparten el valor subjetivo, el razonamiento en el margen y el equilibrio, y discrepan en el método: palabras o matemáticas, equilibrio parcial o general, estado de reposo o proceso.
\end{itemize}

% --------------------------------------------------------------------------
\subsection{Caracterización general de la escuela neoclásica según \autor{Schumpeter}}\label{sec:schumpeter}

\begin{itemize}
\item \autor{Schumpeter} reconoce una escuela de pensamiento por 3 rasgos: un \emph{líder}, unos \emph{seguidores} y un \emph{cuerpo doctrinal} propio, con su objeto y su método \parencite{Schumpeter1954History}. La escuela neoclásica tiene un cuerpo doctrinal común, pero no un solo líder: son varias escuelas nacionales.
  \begin{itemize}
  \item \emph{Líderes}: los 3 pioneros de la revolución marginalista, \autor{William Stanley Jevons}, \autor{Carl Menger} y \autor{Léon Walras}, y \autor{Alfred Marshall}. A \autor{Walras} y a \autor{Marshall} se les considera los padres de la economía neoclásica, porque son los que juntan la oferta y la demanda en un sistema.
  \item \emph{Seguidores}, por escuelas:
    \begin{itemize}
    \item \emph{Cambridge}, la neoclásica, no la postkeynesiana de los años 30 a 60 \vertema{3.A.4}: \autor{Arthur Cecil Pigou}, sucesor de \autor{Marshall} en su cátedra, y, en Inglaterra, \autor{Francis Ysidro Edgeworth} y \autor{Philip Wicksteed};
    \item \emph{Lausana}: \autor{Vilfredo Pareto}, sucesor de \autor{Walras}, y \autor{Enrico Barone};
    \item \emph{Viena}: \autor{Eugen von Böhm-Bawerk} y \autor{Friedrich von Wieser} y, en la generación siguiente, \autor{Ludwig von Mises} y \autor{Friedrich Hayek};
    \item \emph{fuera de las 3}: \autor{Knut Wicksell}, en Suecia, y \autor{John Bates Clark} e \autor{Irving Fisher}, en Estados Unidos.
    \end{itemize}
  \item \emph{Cuerpo doctrinal}:
    \begin{itemize}
    \item \emph{objeto}: la asignación de recursos escasos entre usos alternativos y la formación de los precios relativos que la guían; la distribución de la renta pasa a ser un problema más de precios, los de los factores. Los clásicos estudiaban el crecimiento y el reparto entre clases: el objeto se estrecha y se vuelve estático. \textcite{Robbins1932Essay} lo resumirá en su definición de la economía: la ciencia que estudia la conducta humana como relación entre fines y medios escasos que tienen usos alternativos. En lo agregado, aceptan la ley de \autor{Say} de los clásicos: con precios y salarios flexibles, los mercados se vacían \vertema{3.A.4};
    \item \emph{método}: el individuo que optimiza en el margen y el equilibrio; las diferencias entre escuelas, en el método (\ref{sec:metodo}).
    \end{itemize}
  \end{itemize}
\end{itemize}

\begin{table}[H]
\centering
\caption{Lo que cambia de los clásicos a los neoclásicos}\label{tab:rasgos}
\small
\begin{tabular}{>{\raggedright\arraybackslash}m{2.4cm}>{\raggedright\arraybackslash}m{5.4cm}>{\raggedright\arraybackslash}m{6.2cm}}
\toprule
\rowcolor{tceeazulpalido}\textbf{Rasgo} & \textbf{Clásicos} \vertema{3.A.2} & \textbf{Neoclásicos}\\
\midrule
Pregunta & El crecimiento a largo plazo y el reparto de la renta & \textcolor{tceeazul}{\textbf{La asignación}} de recursos escasos y los precios relativos\\
\addlinespace[0.3em]
Sujeto & Las clases sociales: trabajadores, capitalistas y terratenientes & \textcolor{tceeazul}{\textbf{El individuo}} que optimiza: el consumidor y la empresa\\
\addlinespace[0.3em]
Valor & Objetivo: el coste de producción, en último término el trabajo & \textcolor{tceeazul}{\textbf{Subjetivo}}: la utilidad marginal y, con \autor{Walras} y \autor{Marshall}, la oferta y la demanda a la vez\\
\addlinespace[0.3em]
Distribución & Entre clases; el beneficio, como residuo & \textcolor{tceeazul}{\textbf{Entre factores}}: cada uno cobra su productividad marginal\\
\addlinespace[0.3em]
Método & Verbal y deductivo; economía política, con las instituciones dentro & \textcolor{tceeazul}{\textbf{Análisis marginal}} y equilibrio, con matemáticas en Lausana y, en notas, en Cambridge; sin ellas en Viena\\
\addlinespace[0.3em]
Política & \emph{Laissez-faire} y libre cambio & En general, también; pero el Estado puede corregir los fallos del mercado (\autor{Pigou})\\
\bottomrule
\end{tabular}
\fuente{elaboración propia a partir de \textcite[caps.~8-10]{LandrethColander2006Historia} y \textcite{Robbins1932Essay}. En azul, lo que cambia.}
\end{table}

% --------------------------------------------------------------------------
\subsection{Autores y rango histórico}\label{sec:autores}

\begin{itemize}
\item El marginalismo tiene precursores y una fecha de nacimiento, la de la revolución de 1871-1874; las 3 etapas de su desarrollo, presentadas en la introducción, ordenan los bloques~\ref{sec:revolucion} y~\ref{sec:segunda}.
\end{itemize}

\subsubsection{Precursores del marginalismo}\label{sec:precursores}

\begin{itemize}
\item \emph{De los clásicos} toman los \emph{rendimientos decrecientes} de \autor{Turgot}, la \emph{renta diferencial} de \textcite{Ricardo1817Principles}, que ya es un razonamiento en el margen aplicado a la tierra, el método deductivo y unos mercados en libre competencia.
\item \emph{Aproximaciones pioneras a la utilidad marginal}: antes de 1871, varios autores razonan ya con la \concepto{utilidad}, la capacidad de los bienes de satisfacer necesidades, aunque de forma incompleta o sin sacarle partido para la teoría de los precios:\footnote{\autor{Jeremy Bentham} (1789, 1802) ya había dicho que la utilidad de la riqueza crece cada vez menos, y \textcite{Lloyd1833Lecture} distingue la utilidad total de la marginal, aunque ninguno la conecta con la demanda \parencite{Fonseca2026Phases}.}
  \begin{itemize}
  \item \autor{Daniel Bernoulli} \parencite*{Bernoulli1738Specimen} usa una utilidad marginal decreciente del dinero para resolver la paradoja de San Petersburgo: nadie pagaría una fortuna por un juego cuya ganancia esperada es infinita \vertema{3.A.10}.
  \item \autor{Antoine Augustin Cournot} \parencite*{Cournot1838Recherches} escribe la demanda como una función del precio, $D=F(p)$, y dibuja la demanda y la oferta en el plano precio-cantidad; \autor{Karl Heinrich Rau} (1841) y \autor{Fleeming Jenkin} \parencite*{Jenkin1870Graphic} dibujan ya el cruce de las 2 curvas \parencite{Humphrey1992Marshallian}.\footnote{\autor{Cournot} no funda la demanda en la utilidad: la construye a partir de relaciones observadas entre precios y cantidades. Que la cantidad demandada dependa también de la renta o de otros precios no invalida su análisis, que los toma como dados.}
  \item \autor{Jules Dupuit} \parencite*{Dupuit1844Mesure}, ingeniero de puentes, relaciona la demanda con la utilidad del consumidor y mide con el \emph{excedente del consumidor} lo que vale una obra pública \vertema{3.A.16}.
  \item \autor{Johann Heinrich von Thünen} formula en 1850, en el 2.º volumen de \emph{Der isolirte Staat}, el pago de los factores según su productividad marginal \parencite{Thunen1826Staat}: es el germen de la teoría de la distribución de la segunda generación (\ref{sec:productividad}).
  \item \autor{Hermann Heinrich Gossen} \parencite*{Gossen1854Entwickelung} formula sus 2 leyes:
    \begin{enumerate}
    \item \emph{Ley de la saturación de las necesidades}: la utilidad de cada unidad adicional de un bien disminuye conforme aumenta la cantidad disponible, hasta llegar a la saciedad (utilidad marginal decreciente). Es lo que dice el refrán: «lo poco agrada y lo mucho enfada».
    \item \emph{Ley de la igualdad de las utilidades marginales ponderadas}: el individuo reparte su renta de modo que la última unidad monetaria gastada en cada bien le reporte la misma utilidad, con $U_i$ la utilidad marginal del bien $i$ y $\lambda$ la de la renta \vertema{3.A.8}:
    \funcionclave[eq:gossen]{\frac{U_1}{p_1}=\frac{U_2}{p_2}=\dots=\frac{U_n}{p_n}=\lambda}
    \end{enumerate}
    Algunos manuales añaden una tercera, la \emph{ley de la escasez} (un bien solo tiene valor si es escaso) \parencite{Fonseca2026Phases}, que \autor{Gossen} no formula como tal.\footnote{El libro de \autor{Gossen} pasó inadvertido. \autor{Jevons} lo descubrió años después y le reconoció la prioridad en el prefacio de la 2.ª edición de su \emph{Theory} \parencite{Jevons1879Theory}.}
  \end{itemize}
\end{itemize}

\begin{anotaciones}
\textbf{Más precursores: el valor por la utilidad y la escasez antes de 1871}
\begin{itemize}
\item \emph{La escuela de Salamanca}: en el siglo~XVI, teólogos como \autor{Diego de Covarrubias} y \autor{Martín de Azpilcueta} explican el precio justo por la estimación común de los compradores y por la abundancia o escasez del bien, no por su coste, y \autor{Azpilcueta} relaciona el valor del dinero con su cantidad \parencite{GriceHutchinson1952Salamanca}. La escuela austriaca los reivindicará como antecedentes de su valor subjetivo.
\item \emph{La tradición francesa e italiana}: \textcite{Galiani1751Moneta} y \autor{Étienne Bonnot de Condillac} (1776) ya hacen depender el valor de la utilidad y de la rareza; los clásicos británicos dejaron esa línea en segundo plano \parencite{Fonseca2026Phases}.
\item La idea, por tanto, no era nueva: lo nuevo en 1871-1874 es razonar con la \emph{última} unidad y convertirlo en el centro de una teoría de los precios.
\end{itemize}
\end{anotaciones}

\subsubsection{La revolución marginalista (1871-1874): ideas comunes}\label{sec:ideas-comunes}

\begin{itemize}
\item \textcite{LandrethColander2006Historia} llaman \emph{primera generación} de marginalistas a \autor{Jevons}, \autor{Menger} y \autor{Walras}: sus obras son casi simultáneas e independientes, y comparten 5 ideas:
  \begin{enumerate}
  \item \emph{el valor es subjetivo}: depende de la importancia que cada individuo da a un bien para satisfacer sus necesidades, no del trabajo que lleva incorporado;
  \item \emph{el valor se decide en el margen}: lo fija la utilidad de la última unidad, que decrece con la cantidad y depende, por tanto, de la escasez, y no la utilidad total;
  \item \emph{el principio equimarginal}: el consumidor reparte su renta hasta igualar las utilidades marginales ponderadas por los precios, la segunda ley de \autor{Gossen} \eqref{eq:gossen};
  \item \emph{la utilidad es cardinal} para cada individuo, una magnitud que se puede sumar, pero no se compara entre personas;\footnote{\autor{Jevons} advierte que no hay un denominador común para comparar lo que sienten 2 personas \parencite[cap.~I]{Jevons1871Theory}. \autor{Menger} ni siquiera mide la utilidad: ordena las necesidades por su importancia, y sus números son solo ilustrativos.}
  \item \emph{la oferta queda en segundo plano}: el coste importa solo porque limita la cantidad que llega al mercado, y el valor de los factores se deriva del de los bienes que producen, y no al revés (\ref{sec:menger}).
  \end{enumerate}
\item Los 3 usan el margen de forma distinta:
  \begin{itemize}
  \item \autor{Jevons} y \autor{Menger} buscan cadenas causales sencillas: \autor{Jevons}, en el hogar; \autor{Menger}, en el hogar y en la producción;
  \item \autor{Walras} ve la interdependencia de todas las variables: para él, el margen es solo el primer peldaño de un sistema de equilibrio general.
  \end{itemize}
\item Faltaba juntar las 2 hojas de la tijera. \autor{Marshall} lo hace mercado a mercado, en equilibrio parcial, consciente de la interdependencia pero dejándola en segundo plano, y \autor{Walras}, con todos los mercados a la vez. Con ellos, el precio lo determinan a la vez la oferta y la demanda, y se acaban el valor trabajo, el valor como coste de producción y la distribución como residuo de los clásicos.
\end{itemize}

% --------------------------------------------------------------------------
\subsection{Método}\label{sec:metodo}

\begin{itemize}
\item Una vez vistos los autores y sus fechas, pasamos a lo que los une y los separa: el método.
\item \emph{Lo común}:
  \begin{itemize}
  \item \concepto{individualismo metodológico}: los fenómenos sociales se explican por las decisiones de los individuos, no por las de las clases o los países;
  \item \concepto{análisis marginal}: se decide comparando lo que se gana y lo que se pierde con una unidad más, y el problema económico se plantea como un problema de optimización;
  \item \concepto{equilibrio}: el estado en que nadie quiere cambiar su decisión, al que tienden los mercados y con el que se explican los precios;
  \item \concepto{abstracción}: modelos muy simplificados de hogares que maximizan la utilidad y empresas que maximizan el beneficio, con el método deductivo de los clásicos y el cálculo diferencial: en el texto en Lausana y, en notas, en Cambridge; sin él en Viena.
  \end{itemize}
\item \emph{Lo que los separa}: 3 preguntas de método, que contestan de forma distinta las 3 escuelas (tabla~\ref{tab:escuelas}):
  \begin{itemize}
  \item ¿palabras o matemáticas?
  \item ¿un mercado o todos a la vez?
  \item ¿un estado de reposo o un proceso en el tiempo?
  \end{itemize}
\end{itemize}

\begin{table}[H]
\centering
\caption{Las 3 escuelas neoclásicas}\label{tab:escuelas}
\small
\begin{tabular}{>{\raggedright\arraybackslash}m{2cm}>{\raggedright\arraybackslash}m{4.1cm}>{\raggedright\arraybackslash}m{4.1cm}>{\raggedright\arraybackslash}m{4.1cm}}
\toprule
\rowcolor{tceeazulpalido} & \textbf{Cambridge} & \textbf{Lausana} & \textbf{Viena}\\
\midrule
Fundador & \autor{Marshall} (1890) & \autor{Walras} (1874) & \autor{Menger} (1871)\\
\addlinespace[0.3em]
Seguidores & \autor{Pigou}; en Inglaterra, \autor{Edgeworth} y \autor{Wicksteed} & \autor{Pareto} y \autor{Barone} & \autor{Böhm-Bawerk} y \autor{Wieser}; después, \autor{Mises} y \autor{Hayek}\\
\addlinespace[0.3em]
Método & \textcolor{tceeazul}{\textbf{Verbal}}, con las matemáticas en notas y apéndices; cerca de los hechos & \textcolor{tceeazul}{\textbf{Matemático}}: la economía pura, como la mecánica & \textcolor{tceeazul}{\textbf{Verbal y deductivo}}, a partir de la acción individual\\
\addlinespace[0.3em]
Equilibrio & \textcolor{tceeazul}{\textbf{Parcial}}: un mercado, \emph{ceteris paribus}, con plazos & \textcolor{tceeazul}{\textbf{General}}: todos los mercados a la vez & \textcolor{tceeazul}{\textbf{Proceso}}: el mercado como descubrimiento, no como estado de reposo\\
\addlinespace[0.3em]
Valor & Utilidad y coste a la vez: las tijeras & \emph{Rareté} y oferta y demanda de todos los bienes & Subjetivo; los factores, por imputación\\
\addlinespace[0.3em]
Capital e interés & Cuasirrentas; la espera como coste & El capital, dentro del sistema de ecuaciones & Rodeos de producción y preferencia por el presente (\autor{Böhm-Bawerk})\\
\addlinespace[0.3em]
Dinero & Saldos de caja: $M^D=k\,P\,Y$ & El dinero, en la última etapa del sistema (la circulación) & Origen del dinero (\autor{Menger}); dinero y ciclo (\autor{Mises}, \autor{Hayek})\\
\addlinespace[0.3em]
Bienestar & Excedente del consumidor; coste social y \newline impuesto (\autor{Pigou}) & Máximo de satisfacción; óptimo de \autor{Pareto} & Crítica del cálculo económico socialista\\
\addlinespace[0.3em]
En la exposición & \ref{sec:marshall} & \ref{sec:walras} & \ref{sec:menger}\\
\midrule
\multicolumn{4}{>{\centering\arraybackslash}m{15.5cm}}{\cellcolor{tceesalvia}Fuera de las 3: \autor{Wicksell}, en Suecia, y \autor{Clark} y \autor{Fisher}, en Estados Unidos (\ref{sec:segunda})}\\
\bottomrule
\end{tabular}
\fuente{elaboración propia a partir de \textcite[caps.~8-13]{LandrethColander2006Historia} y \textcite{Fonseca2026Neoclassicals}. En azul, el rasgo de método que distingue a cada escuela.}
\end{table}

\begin{itemize}
\item \emph{La disputa sobre el método}: en 1883, \autor{Menger} defiende la teoría abstracta frente a la escuela histórica alemana en sus \emph{Untersuchungen} \parencite{Menger1883Untersuchungen}; \autor{Gustav von Schmoller}, que quería una economía inductiva, construida desde la historia y las instituciones, responde con una reseña durísima \parencite{Schmoller1883Methodologie}, y \autor{Menger} replica con \emph{Die Irrthümer des Historismus} \parencite{Menger1884Irrthumer}. Es la \concepto{disputa sobre el método} (\emph{Methodenstreit}) \vertema{3.A.1}.
  \begin{itemize}
  \item «Cada loco con su tema»: cada bando defendía su método como el único científico, y la disputa enconó durante años la economía en lengua alemana.
  \item Los austriacos quedan como \emph{deductivos}: parten de unos pocos principios de la acción humana, sin matemáticas, y en eso se distinguen tanto de los historicistas como de \autor{Walras}.
  \end{itemize}
\item \emph{Recapitulación}: la escuela neoclásica tiene varios líderes y 3 escuelas que comparten el valor subjetivo, el margen y el equilibrio, y discrepan en el método. Estudiado el conjunto, pasamos a sus fundadores: los 3 de la revolución marginalista y \autor{Marshall}, que la convierte en sistema.
\end{itemize}

% ==========================================================================
\section{Revolución marginalista (1871-1874)}\label{sec:revolucion}
\minutos{12,5}

\begin{itemize}
\item Una vez caracterizada la escuela, pasamos a la 1.ª etapa de \autor{Fonseca}, la revolución marginalista, con sus 4 fundadores, cada uno con la escuela que creó o en la que se apoyó: los 3 de 1871-1874 y \autor{Marshall}, que la consolida en 1890.
\item \emph{Conclusión anticipada}: entre 1871 y 1874 el valor pasa a depender de la utilidad marginal; \autor{Menger} lo lleva a los factores y al tiempo, \autor{Walras} a todos los mercados a la vez y, en 1890, \autor{Marshall} une la utilidad y el coste en las tijeras de la oferta y la demanda.
\end{itemize}

% --------------------------------------------------------------------------
\subsection{\autor{William Stanley Jevons} (1835-1882)}\label{sec:jevons}

\paragraph{a) Contexto, método y obra}

\begin{itemize}
\item \emph{Contexto}: químico y ensayador en la casa de la moneda de Sídney, lógico y estadístico, profesor en Manchester y en Londres. Antes de la \emph{Theory} se había hecho famoso con \emph{The Coal Question} \parencite{Jevons1865Coal}.
\item \emph{Método}: matemático y utilitarista. Para \autor{Jevons}, la economía tiene que ser una ciencia matemática, porque trata de cantidades, y es un cálculo del placer y del dolor a la manera de \autor{Bentham} \parencite[cap.~I]{Jevons1871Theory}.
\end{itemize}

\paragraph{b) La utilidad y el intercambio}

\begin{itemize}
\item \emph{La utilidad}: es \emph{aditiva}, la suma de las que dan los bienes por separado, y lo que cuenta es el \concepto{grado final de utilidad}, la utilidad marginal, que decrece:
  \begin{equation}
  U(\vect{x})=u_1(x_1)+u_2(x_2)+\dots+u_n(x_n),\qquad U_i=u_i'(x_i),\quad u_i''(x_i)<0. \label{eq:jevons}
  \end{equation}
  Con esa forma, la utilidad marginal de un bien no depende de cuánto se tenga de los demás: no hay sustitutivos ni complementarios. El consumidor reparte su renta según la segunda ley de \autor{Gossen} \eqref{eq:gossen}.
\item \emph{La paradoja del agua y los diamantes}, resuelta: el agua da mucha más utilidad total, pero hay tanta que la última unidad vale poco; de los diamantes hay tan pocos que la última vale mucho, y es la marginal la que fija el precio (gráfico~\ref{fig:agua-diamantes}). Volvemos al refrán del principio: nadie sabe lo que vale el agua hasta que el pozo se seca, porque lo que vale es la última unidad, y su valor depende de cuánta agua quede. La escasez que cuenta es la que se desea, no la mera rareza (anexo~\ref{anexo:escasez}).
\end{itemize}

\begin{figure}[H]
\centering
\caption{La paradoja del agua y los diamantes: utilidad total y utilidad marginal}\label{fig:agua-diamantes}
\grafico[0.9\textwidth]{agua-diamantes}
\fuente{elaboración propia a partir de \textcite[libro~I, cap.~IV]{Smith1776Wealth} y \textcite[cap.~III]{Menger1871Grundsatze}. El área bajo cada curva hasta la cantidad disponible es la utilidad total; la altura en ese punto, la marginal, que es la que fija el precio. Curvas ilustrativas.}
\end{figure}

\begin{itemize}
\item \emph{La ley de la indiferencia}: «en un mismo mercado abierto, en un momento dado, no puede haber 2 precios para la misma clase de artículo» \parencite[págs.~91-92; traducción propia]{Jevons1871Theory}. Es la \concepto{ley del precio único}, que le permite escribir la relación de intercambio como un cociente de precios.
\item \emph{El intercambio}: 2 personas, $A$ con el bien 1 y $B$ con el bien 2, cambian hasta que la razón de intercambio, $x_2/x_1$, iguala el cociente de las utilidades marginales de cada una (anexo~\ref{anexo:derivaciones}):
  \begin{equation}
  \mathrm{RMS}^A_{12}=\frac{U^A_1}{U^A_2}=\frac{x_2}{x_1}=\frac{U^B_1}{U^B_2}=\mathrm{RMS}^B_{12}. \label{eq:intercambio-jevons}
  \end{equation}
  \autor{Edgeworth} mostrará que, con solo 2 personas, el resultado está indeterminado (\ref{sec:interes}).
\item \emph{La cadena de \autor{Jevons}}: «el coste de producción determina la oferta; la oferta determina el grado final de utilidad; el grado final de utilidad determina el valor» \parencite[traducción propia]{Jevons1879Theory}; en la 1.ª edición lo decía en prosa \parencite[págs.~159-160]{Jevons1871Theory}. El coste importa, pero solo a través de la cantidad que llega al mercado.
\item \emph{El trabajo}: se trabaja hasta que la penosidad de la última hora iguala la utilidad de lo que con ella se obtiene, el antecedente de la oferta de trabajo neoclásica \vertema{3.A.25}.
\end{itemize}

\paragraph{c) Otras aportaciones}

\begin{itemize}
\item \emph{Salarios y beneficios}: niega que se muevan en sentido contrario; con un producto variable, esa doctrina es «radicalmente falaz», porque intenta despejar 2 incógnitas con una sola ecuación, y la teoría del fondo de salarios es «del todo ilusoria» \parencite[págs.~256-258; traducción propia]{Jevons1871Theory}.
\item \emph{La paradoja de \autor{Jevons}}: usar el carbón con más eficiencia no reduce su consumo, lo aumenta, porque abarata su uso: «es una completa confusión de ideas suponer que el uso económico del combustible equivale a un menor consumo. Lo cierto es lo contrario» \parencite[pág.~103; traducción propia]{Jevons1865Coal}. Es el \emph{efecto rebote} de la economía de la energía.
\item \emph{El ciclo}: atribuye las crisis comerciales a las manchas solares, que alterarían las cosechas y el precio del grano (memorias de 1875 y 1878, recogidas en \textcite{Jevons1884Investigations}; \cite{HET2026Jevons}): una teoría \emph{real} del ciclo. Los \emph{equilibrios de manchas solares} de \textcite{Azariadis1981Self} y \textcite{CassShell1983Sunspots} toman de ahí el nombre para algo distinto: una variable sin relación con los fundamentos que coordina las expectativas y elige entre varios equilibrios \vertema{3.A.42}.
\end{itemize}

\paragraph{d) Valoración}

\begin{itemize}
\item \autor{Jevons} da la formulación matemática más clara de la utilidad marginal, pero su utilidad aditiva y su desinterés por la oferta dejan la teoría a medias: no explica qué determina el coste. Es lo que intentan, por caminos opuestos, \autor{Menger} y sus discípulos.
\end{itemize}

% --------------------------------------------------------------------------
\subsection{\autor{Carl Menger} (1840-1921) y la escuela de Viena}\label{sec:menger}

\paragraph{a) Contexto, método y obra}

\begin{itemize}
\item \emph{Contexto}: periodista económico y funcionario en Viena, publica los \emph{Grundsätze} en 1871 y en 1873 entra en la Universidad de Viena, donde forma la escuela austriaca.
\item \emph{Método}: deductivo y verbal, sin matemáticas, a partir de las necesidades de los individuos y de lo que saben sobre cómo satisfacerlas (lo que después \autor{Mises} llamará praxeología). Lo defiende frente a la escuela histórica en la disputa sobre el método (\ref{sec:metodo}).
\end{itemize}

\paragraph{b) La teoría del valor}

\begin{itemize}
\item \emph{Los bienes}: una cosa es un \emph{bien} si satisface una necesidad y se sabe cómo usarla, y es un \emph{bien económico} si hay menos de la que se necesita.
  \begin{itemize}
  \item Los \concepto{bienes de orden inferior} (o de primer orden) satisfacen las necesidades directamente: el pan.
  \item Los \concepto{bienes de orden superior} lo hacen de forma indirecta, porque sirven para producir los anteriores: la harina, el horno, el trabajo del panadero. Son \emph{complementarios}: ninguno sirve sin los demás.
  \item Su valor se \concepto{imputa} desde el de los bienes de consumo que ayudan a producir: no son los costes los que fijan el precio, sino al revés \parencite{Menger1871Grundsatze}. Si nadie quisiera pan, el horno no valdría nada, por mucho que hubiera costado.
  \end{itemize}
\item \emph{El valor}: cada individuo ordena sus necesidades por su importancia y la de las unidades sucesivas de cada bien, que es decreciente, y reparte lo que tiene hasta igualar la importancia de la última necesidad que cubre en cada uso (tabla~\ref{tab:menger}).
  \begin{itemize}
  \item El \emph{valor de una unidad} es la importancia de la necesidad menos urgente que depende de ella: lo que se perdería si se perdiera esa unidad. Es la utilidad marginal, sin la palabra, y es otra vez el refrán del enganche: con el pozo casi seco, el último vaso cubre la sed, y con el pozo lleno, regar las flores.
  \item La escasez que importa es relativa al deseo, no a la cantidad física: los huevos podridos, aunque haya pocos, no son escasos, porque nadie los quiere \parencite[cap.~III]{Robbins1932Essay}.
  \end{itemize}
\item \emph{El origen del dinero}: el dinero no lo inventó el Estado. Cada uno acepta a cambio el bien más fácil de vender, aunque no lo quiera consumir, porque le ahorra buscar a quien quiera justo lo que ofrece; poco a poco, ese bien se convierte en el medio de cambio de todos \parencite{Menger1892Origin}. Es el ejemplo de una institución que nadie ha diseñado \vertema{3.A.35}.
\end{itemize}

\begin{table}[H]
\centering
\caption{La escala de necesidades de \autor{Menger}: importancia de cada unidad sucesiva en 5 clases de necesidades}\label{tab:menger}
\small
\begin{tabular}{>{\centering\arraybackslash}m{2cm}*{5}{>{\centering\arraybackslash}m{1.3cm}}}
\toprule
\rowcolor{tceeazulpalido}\textbf{Unidad} & \textbf{I} & \textbf{II} & \textbf{III} & \textbf{IV} & \textbf{V}\\
\midrule
1.ª & \textcolor{tceeverde}{\textbf{10}} & \textcolor{tceeverde}{\textbf{9}} & \textcolor{tceeverde}{\textbf{8}} & 7 & 6\\
2.ª & \textcolor{tceeverde}{\textbf{9}} & \textcolor{tceeverde}{\textbf{8}} & 7 & 6 & 5\\
3.ª & \textcolor{tceeverde}{\textbf{8}} & 7 & 6 & 5 & 4\\
4.ª & 7 & 6 & 5 & 4 & 3\\
5.ª & 6 & 5 & 4 & 3 & 2\\
\midrule
\multicolumn{6}{>{\centering\arraybackslash}m{10.6cm}}{\cellcolor{tceesalvia}Con 6 unidades: 3 a la clase I, 2 a la II y 1 a la III. La última unidad de cada clase cubre una necesidad de importancia 8, y 8 es el valor de cada unidad}\\
\bottomrule
\end{tabular}
\fuente{elaboración propia a partir de \textcite[cap.~III]{Menger1871Grundsatze}, que usa 10 clases (la I, la más urgente, como la alimentación) y una escala de 10 a 0. En verde, las necesidades que se cubren.}
\end{table}


\paragraph{c) La escuela de Viena: capital, interés, dinero y ciclo}

\subparagraph{i) \autor{Eugen von Böhm-Bawerk} (1851-1914): el capital y el interés}

\begin{itemize}
\item ¿Por qué se paga un interés? En el primer volumen de \emph{Kapital und Kapitalzins} \parencite{BohmBawerk1884Geschichte} rechaza las respuestas anteriores (la productividad del capital, la abstinencia de \autor{Senior} y la explotación de \autor{Marx}), y en el segundo \parencite{BohmBawerk1889Positive} da la suya: los bienes presentes valen más que los futuros, y el interés es esa prima, el \emph{agio}. Lo dice el refrán: «más vale un toma que dos te daré».
\item Las causas son 3:
  \begin{enumerate}
  \item se espera estar mejor provisto en el futuro, y lo que se tendrá entonces se valora menos;
  \item se infravalora el futuro, por falta de imaginación o de voluntad o por la brevedad de la vida;
  \item los bienes presentes son \emph{técnicamente superiores}: permiten métodos de producción más largos, los \concepto{rodeos de producción}, que rinden más (pescar con red en vez de con la mano), con un mayor \emph{periodo medio de producción}.
  \end{enumerate}
\item Las 2 primeras son subjetivas, la preferencia temporal; la 3.ª es técnica. \autor{Fisher} aceptará las 2 primeras y rechazará la 3.ª (\ref{sec:interes}); el periodo medio de producción no se puede medir cuando hay muchos bienes de capital distintos, un anticipo de la controversia del capital \vertema{3.A.4}.
\item En \emph{Zum Abschluss des Marxschen Systems} \parencite{BohmBawerk1896Abschluss} critica la teoría del valor trabajo y de la explotación de \autor{Marx}: el valor del libro I de \emph{El capital} no es compatible con los precios de producción del libro III \vertema{3.A.2}.
\end{itemize}

\subparagraph{ii) \autor{Friedrich von Wieser} (1851-1926): la imputación y el coste de oportunidad}

\begin{itemize}
\item Acuña el término \emph{Grenznutzen}, utilidad marginal \parencite{Wieser1884Ursprung}, y hace de la imputación una teoría general: a cada factor complementario se le imputa su contribución al valor del producto \parencite{Wieser1889Naturliche,HET2026Wieser}.
\item Hace del \emph{coste alternativo}, lo que se deja de producir con los mismos factores, la base del coste \parencite{Wieser1889Naturliche}; el nombre inglés, \emph{opportunity cost}, es de \textcite{Green1894Pain}.
\end{itemize}

\subparagraph{iii) \autor{Ludwig von Mises} (1881-1973) y \autor{Friedrich Hayek} (1899-1992): dinero, ciclo y conocimiento}

\begin{itemize}
\item \emph{El dinero}: \textcite{Mises1912Theorie} aplica al dinero la utilidad marginal: hoy se demanda dinero por el poder adquisitivo que tuvo ayer, y así hacia atrás, hasta el día en que era una mercancía que se usaba por sí misma.
\item \emph{La teoría austriaca del ciclo} \parencite{Mises1912Theorie,Hayek1931Prices}:
  \begin{itemize}
  \item si los bancos prestan por debajo del tipo natural de \autor{Wicksell} (\ref{sec:interes}), los empresarios emprenden procesos de producción más largos de lo que el ahorro voluntario puede sostener;
  \item el auge dura lo que dura el crédito; cuando se frena, esos proyectos se abandonan y llega la crisis, que liquida los errores. Lo que hay que evitar es el auge \vertema{3.A.42}.
  \end{itemize}
\item \emph{El cálculo económico en el socialismo}: sin propiedad ni mercados de los bienes de capital no hay precios con los que calcular \parencite{Mises1920Wirtschaftsrechnung}. \textcite{Lange1936Socialism} responde que un planificador podría tantear los precios como el mercado de \autor{Walras} (\ref{sec:walras}).
\item \emph{El mercado como proceso}: el problema económico no es resolver ecuaciones, sino aprovechar un conocimiento disperso entre millones de personas, que los precios resumen \parencite{Hayek1945Use}; la competencia es un proceso de descubrimiento en el tiempo, no el estado de reposo de la competencia perfecta \parencite{Hayek1948Meaning} \vertema{3.A.16}. En esa línea, \textcite{Schumpeter1912Theorie} pone en el centro al empresario que innova y rompe el equilibrio, y \textcite{Kirzner1973Competition}, al que descubre oportunidades.\footnote{\autor{Hayek} recibe en 1974, con \autor{Gunnar Myrdal}, el Premio Nobel por sus trabajos sobre el dinero y las fluctuaciones y por su análisis de la interdependencia de los fenómenos económicos, sociales e institucionales \parencite{Nobel1974}.}
\end{itemize}

\paragraph{d) Valoración}

\begin{itemize}
\item Viena da el subjetivismo más coherente y lo lleva al tiempo, al capital, al dinero y al conocimiento, pero sin matemáticas ni un sistema cerrado. El sistema completo de precios y cantidades lo construye, con matemáticas, Lausana.
\end{itemize}

% --------------------------------------------------------------------------
\subsection{\autor{Léon Walras} (1834-1910) y la escuela de Lausana}\label{sec:walras}

\paragraph{Contexto, método y obra}

\begin{itemize}
\item \emph{Contexto}: hijo del economista \autor{Auguste Walras}, de quien toma la idea de la \emph{rareté}, e influido por \autor{Cournot}, compañero de su padre en la Escuela Normal, enseña en Lausana desde 1870. Los \emph{Éléments} salen en 2 entregas, en 1874 y 1877 \parencite{Walras1874Elements}.\footnote{\autor{Walras}, padre de una teoría que hoy algunos tachan de conservadora, era socialista: defendía nacionalizar la tierra y los monopolios \parencite{Daniel2021Histoire}.}
\item \emph{Método}: distingue la economía \emph{pura}, la teoría del valor de cambio, una ciencia como la mecánica que debe escribirse con matemáticas; la \emph{aplicada}, la de la producción, y la \emph{social}, la de la justicia en el reparto.
\end{itemize}

\paragraph{El equilibrio general}

\subparagraph{a) Idea}

\begin{itemize}
\item Todo depende de todo: el precio del pan depende del de la harina, del salario del panadero y de lo que los hogares gastan en otras cosas. \autor{Walras} escribe la economía entera como un sistema de ecuaciones y se pregunta si hay unos precios que vacíen todos los mercados a la vez \vertema{3.A.21}.
\end{itemize}

\subparagraph{b) Modelo}

\begin{itemize}
\item \emph{i) Supuestos}: competencia en todos los mercados; consumidores que maximizan su utilidad y venden los servicios de sus factores; empresas que producen con coeficientes de fabricación, y un bien que sirve de \concepto{numerario}, en el que se miden los demás precios.
\item \emph{ii) Desarrollo}: el sistema crece en 4 etapas: el intercambio, la producción, la capitalización (el ahorro y los bienes de capital) y la circulación (el dinero).
  \begin{itemize}
  \item La \emph{rareté}, la intensidad de la última necesidad satisfecha, es la utilidad marginal: en el equilibrio del consumidor, las \emph{raretés} son proporcionales a los precios \eqref{eq:gossen}, y de ahí \autor{Walras} deriva las demandas como funciones de todos los precios. Es la versión general de la demanda recíproca de \autor{Mill}.
  \item Como cada agente respeta su presupuesto, el valor de los excesos de demanda de todos los mercados suma cero a cualquier precio: es la \concepto{ley de \autor{Walras}}. Si todos los mercados menos uno están en equilibrio, también lo está el último, y solo se determinan los precios relativos.
  \item \autor{Walras} cuenta ecuaciones e incógnitas, las ve iguales y concluye que el sistema está determinado.
  \item Para explicar cómo se llega a esos precios recurre al \concepto{tanteo} (\emph{tâtonnement}): el precio sube donde hay exceso de demanda y baja donde hay exceso de oferta, hasta vaciar todos los mercados.\footnote{\autor{Walras} no habla de un subastador: escribe que en el mercado se vocean precios (\emph{on y crie}) \parencite[§~126, pág.~127]{Walras1874Elements}. Tras las críticas de \autor{Bertrand} (1883) y \autor{Edgeworth}, en la 4.ª edición (1900) el tanteo pasa a hacerse con compromisos escritos (\emph{bons}) que solo se ejecutan en el equilibrio \parencite{Walker1987Tatonnement} \vertema{3.A.21}.}
  \end{itemize}
\item \emph{iii) Implicaciones}: precios y cantidades se determinan a la vez, sin una cadena causal como la de \autor{Jevons}; en el equilibrio el empresario no gana ni pierde, y el precio de un bien de capital es el valor capitalizado de los servicios que presta.
\end{itemize}

\subparagraph{c) Valoración}

\begin{itemize}
\item Contar ecuaciones no basta: pueden no tener una solución con precios y cantidades no negativos. La primera demostración de que el equilibrio existe es de \textcite{Wald1936Gleichungssysteme}, y la general, de \textcite{ArrowDebreu1954Existence}. Que el tanteo converja tampoco está garantizado \vertema{3.A.21}.
\item Su obra tardó en leerse fuera del mundo francófono: la traduce al inglés \autor{Jaffé} en 1954 \parencite{Walras1954Elements}, el año de la demostración de \autor{Arrow} y \autor{Debreu}. Desde entonces es el marco de la microeconomía y de buena parte de la macroeconomía.
\end{itemize}

\paragraph{La escuela de Lausana: \autor{Pareto} y \autor{Barone}}

\begin{itemize}
\item \autor{Vilfredo Pareto} (1848-1923), ingeniero, sucede a \autor{Walras} en Lausana en 1893 \parencite{HET2026Pareto}. En el \emph{Cours} \parencite{Pareto1896Cours} y el \emph{Manuale} \parencite{Pareto1906Manuale} sustituye la utilidad por la \emph{ofelimidad}, que solo ordena las cestas: es la utilidad ordinal \vertema{3.A.8}. Su criterio de óptimo, con el bienestar (\ref{sec:interes}).
\item Su discípulo \autor{Enrico Barone} \parencite*{Barone1908Ministro} muestra que un ministerio de la producción que resolviera las mismas ecuaciones llegaría al mismo resultado que la competencia: el punto de partida del debate sobre el cálculo económico en el socialismo.
\end{itemize}

% --------------------------------------------------------------------------
\subsection{\autor{Alfred Marshall} (1842-1924) y la escuela de Cambridge}\label{sec:marshall}

\paragraph{Contexto, método y obra}

\begin{itemize}
\item \emph{Contexto}: matemático de formación, llega a la economía desde la ética, preocupado por la pobreza. Es catedrático en Cambridge de 1885 a 1908, y sus \emph{Principles of Economics} \parencite{Marshall1890Principles}, con 8 ediciones hasta 1920, sustituyen a los de \textcite{Mill1848Principles} como el manual de varias generaciones.
\item \emph{Método}:
  \begin{itemize}
  \item la economía es «el estudio de la humanidad en los asuntos corrientes de la vida» \parencite[libro~I, cap.~I, §~1; traducción propia]{Marshall1890Principles}: una ciencia de hechos, con la historia y las instituciones a la vista;
  \item \emph{equilibrio parcial}: aparta por un tiempo las demás causas en el corral del \emph{ceteris paribus} y estudia un mercado aislado \parencite[libro~V, cap.~V, §~2]{Marshall1890Principles} \vertema{3.A.16}. Reconoce la interdependencia, pero la deja en segundo plano;
  \item \emph{las matemáticas, en notas y apéndices}: sirven para pensar, pero el texto tiene que entenderse sin ellas;\footnote{En una carta de 1906 a \autor{Arthur Bowley} resume sus reglas: usar las matemáticas como un lenguaje abreviado, no como motor de la investigación; traducir el resultado a palabras; ilustrarlo con ejemplos de la vida real y quemar las matemáticas; si no se encuentran ejemplos, quemar también lo traducido \parencite[pág.~427]{Pigou1925Memorials}.}
  \item \emph{continuidad} con los clásicos, de los que se considera heredero, y con la biología: el lema de los \emph{Principles} es \emph{natura non facit saltum}.
  \end{itemize}
\item \emph{Frente a \autor{Walras}}: toma la cantidad como variable independiente y pregunta qué precio se pagaría por ella, el \emph{precio de demanda} $p^D(q)$, por eso pone el precio en el eje vertical; \autor{Walras} trata el precio como la variable independiente \vertema{3.A.9}.
\end{itemize}

\paragraph{El equilibrio parcial y los periodos de mercado}

\subparagraph{a) Idea}

\begin{itemize}
\item ¿Lo fija la utilidad o el coste? Los 2: «tan razonable sería discutir si es la hoja superior o la inferior de unas tijeras la que corta un papel como si el valor lo rige la utilidad o el coste de producción» \parencite[libro~V, cap.~III, §~7; traducción propia]{Marshall1890Principles}. «Una mano lava la otra, y las dos, la cara»: oferta y demanda actúan a la vez.
\item Lo que cambia con el plazo es el peso de cada hoja: cuanto más corto el periodo, más importa la demanda; cuanto más largo, más el coste \parencite[libro~V, cap.~III, §~7]{Marshall1890Principles}.
\end{itemize}

\subparagraph{b) Modelo}

\begin{itemize}
\item \emph{i) Supuestos}: un mercado pequeño respecto al conjunto de la economía, en competencia, con todo lo demás dado; la \emph{utilidad marginal del dinero es constante}, y la \concepto{empresa representativa}, una empresa de tamaño normal con acceso a las economías internas y externas de la industria, resume el coste \parencite[libro~IV, cap.~XIII]{Marshall1890Principles}.\footnote{Con la notación de hoy, la utilidad marginal del dinero constante equivale a una utilidad cuasilineal en el dinero que se gasta en lo demás \vertema{3.A.16}. \autor{Marshall} no la escribió así: es una lectura moderna.}
\item \emph{ii) Desarrollo}:
  \begin{itemize}
  \item \emph{La demanda}: con la utilidad marginal del dinero constante, el precio de demanda de cada unidad es su utilidad marginal medida en dinero, y decrece: es la \emph{ley de la demanda}. \autor{Marshall} define 2 herramientas que siguen en uso:
    \begin{itemize}
    \item la \concepto{elasticidad} de la demanda, $\varepsilon_D=\frac{dq}{dp}\frac{p}{q}$, cuánto cambia en porcentaje la cantidad cuando el precio cambia un 1~\% \parencite[libro~III, cap.~IV]{Marshall1890Principles};
    \item el \concepto{excedente del consumidor} (EC), lo que se estaría dispuesto a pagar por encima de lo que se paga, el área bajo la demanda y sobre el precio, $\mathrm{EC}=\int_0^{q}p^D(s)\,ds-p\,q$ \parencite[libro~III, cap.~VI]{Marshall1890Principles}; con la utilidad marginal del dinero constante es una medida del bienestar, y \autor{Marshall} la usa ya para medir lo que cuesta un impuesto \vertema{3.A.16}.
    \end{itemize}
  \item \autor{Marshall} distingue el \emph{movimiento a lo largo} de la demanda, cuando cambia el precio, de su \emph{desplazamiento}, cuando cambian la renta, los gustos o los precios de otros bienes.
  \item \emph{La oferta}: el precio de oferta $p^S(q)$ es el que cubre el coste de producir $q$. Las \emph{economías internas} dependen del tamaño de la empresa y las \emph{externas}, del de la industria, como las de los distritos industriales \parencite[libro~IV, caps.~IX-XIII]{Marshall1890Principles}.
  \item \emph{Los periodos}, según el tiempo que tienen los vendedores para ajustarse \parencite[libro~V, cap.~V]{Marshall1890Principles}:
    \begin{itemize}
    \item el \concepto{periodo de mercado}: la cantidad está dada y el precio lo fija la demanda;
    \item el \concepto{corto plazo}: se produce más con el equipo existente; lo que obtiene el equipo fijo por encima del coste variable es una \concepto{cuasirrenta}, porque desaparece cuando la capacidad se ajusta \parencite[libro~V, cap.~IX]{Marshall1890Principles};
    \item el \concepto{largo plazo}: cambian el equipo y el número de empresas, y el precio tiende al \concepto{valor normal}, el coste de la empresa representativa;
    \item el \emph{periodo secular}: cambian la técnica, la población y el capital.
    \end{itemize}
  \item \emph{El mercado de pescado} \parencite[libro~V, cap.~V, §~4]{Marshall1890Principles} (gráfico~\ref{fig:marshall-periodos}): la oferta de cada día ($MS$) es vertical, porque el pescado ya está en la lonja; la oferta normal ($NS$), la que se ofrecería con tiempo para ajustar la flota, es creciente.
    \begin{itemize}
    \item Si la demanda aumenta de $D_0$ a $D_1$ (en el ejemplo, porque una epidemia del ganado encarece la carne), el primer día la cantidad no cambia y el precio sube hasta vaciar la lonja: es el \concepto{equilibrio temporal} ($B$), con cuasirrentas para los pescadores.
    \item Con tiempo salen más barcos y el mercado va hacia el \concepto{equilibrio normal} ($C$), con un precio mayor que el inicial ($A$) y menor que el del primer día.
    \item Con rendimientos crecientes en la pesca, \autor{Marshall} admite una oferta normal decreciente \parencite[libro~V, cap.~V, §~4]{Marshall1890Principles}, que hoy es la industria de costes decrecientes \vertema{3.A.16}.
    \end{itemize}
  \item \emph{El ajuste}: la cantidad aumenta si el precio de demanda supera al de oferta \parencite[libro~V, cap.~III, §~6]{Marshall1890Principles}: es el \emph{ajuste marshalliano}, por cantidades, frente al \emph{walrasiano}, por precios \vertema{3.A.16}.
  \end{itemize}
\item \emph{iii) Implicaciones}: el precio vacía la lonja cada día y la cantidad se ajusta después: en este mercado no queda pescado sin vender. Las 2 cosas no se contradicen: dentro de cada día manda el precio, que vacía la lonja, y de un día para otro la cantidad sigue la diferencia entre el precio de demanda y el de oferta, que es el ajuste marshalliano por cantidades. Pero el orden importa: en \autor{Marshall} se ajustan antes los precios que las cantidades; \autor{Keynes} invertirá el orden, en la lectura de \textcite{Leijonhufvud1968Keynesian} \vertema{3.A.4}.
\end{itemize}

\begin{figure}[H]
\centering
\caption{El mercado de pescado de \autor{Marshall}: equilibrio temporal ($B$) y equilibrio normal ($C$) tras un aumento de la demanda}\label{fig:marshall-periodos}
\grafico[0.5\textwidth]{marshall-periodos}
\fuente{elaboración propia a partir de \textcite[fig.~1.1]{DeVroey2016History} y \textcite[libro~V, cap.~V]{Marshall1890Principles}.}
\end{figure}

\subparagraph{c) Valoración}

\begin{itemize}
\item Es la base de la microeconomía de los mercados \vertema{3.A.16}. \textcite{Sraffa1926Laws} le encuentra una contradicción: los rendimientos crecientes internos a la empresa no son compatibles con la competencia, y las economías externas no lo son con el equilibrio parcial, porque afectan a otros mercados. La salida será la competencia imperfecta \vertema{3.A.18}.
\end{itemize}

\begin{anotaciones}
\textbf{Más sobre la demanda y los mercados de \autor{Marshall}}
\begin{itemize}
\item \emph{El efecto renta}: con la utilidad marginal del dinero constante, el efecto renta de un cambio de precio es despreciable y la demanda siempre tiene pendiente negativa. Desde la 3.ª edición (1895), \autor{Marshall} recoge la excepción que atribuye a \autor{Robert Giffen}: cuando sube el pan, las familias pobres compran más pan, porque ya no pueden pagar carne \parencite[libro~III, cap.~VI]{Marshall1890Principles}. Es el \emph{bien Giffen}, que la ecuación de \autor{Slutsky} explicará con el efecto renta \vertema{3.A.8}.
\item \emph{El mercado de trabajo}: los \emph{Principles} no le dedican un análisis propio; el salario sale, como cualquier precio, de la oferta y la demanda. Lo desarrollarán \textcite{Hicks1932Theory}, con el paro friccional y los sindicatos, y \textcite{Pigou1933Theory} \vertema{3.A.25}.
\item \emph{Los precios intervenidos}: con un precio mínimo por encima del de equilibrio, la cantidad intercambiada la fija el lado corto del mercado, la demanda, y queda oferta sin vender.
\end{itemize}
\end{anotaciones}

\paragraph{Distribución, dinero y comercio}

\begin{itemize}
\item \emph{La distribución}: los precios de los factores salen de la oferta y la demanda de cada uno, como los de los bienes. A la tierra, el trabajo y el capital añade un 4.º factor, la \emph{organización}.
  \begin{itemize}
  \item La demanda de cada factor depende de su \emph{producto neto marginal}, lo que añade al valor de la producción la última unidad.
  \item La oferta, de su coste: el de criar y formar al trabajador o el de la \emph{espera} del ahorrador, palabra con la que sustituye la \emph{abstinencia} de \autor{Senior}.
  \item La \emph{renta} es el excedente de un factor de oferta fija, como la tierra; la \emph{cuasirrenta}, la de un equipo fijo a corto plazo \vertema{3.A.25}.
  \end{itemize}
\item \emph{El dinero}: cada individuo desea tener en dinero una fracción $k$ de su renta: es el \concepto{enfoque de los saldos de caja de Cambridge}, $M^D=k\,P\,Y$, con $k=1/V$. Lo expone \autor{Marshall} en declaraciones ante comisiones oficiales y en \textcite{Marshall1923}, y lo formaliza \textcite{Pigou1917Value} \vertema{3.A.35}.
\item \emph{El comercio}: con las \emph{curvas de oferta recíproca}, que dibujan la demanda recíproca de \autor{Mill}, determina la relación real de intercambio; aparecen en \emph{The Pure Theory of Foreign Trade}, que circula en impresión privada desde 1879 \parencite{Marshall1879Pure} \vertema{3.B.5}. Son la herramienta con la que después se estudia el arancel óptimo \vertema{3.B.7}.
\end{itemize}

\paragraph{La escuela de Cambridge: \autor{Pigou}}

\begin{itemize}
\item \autor{Arthur Cecil Pigou} (1877-1959) sucede a \autor{Marshall} en 1908 \parencite{HET2026Pigou} y sistematiza su obra: formaliza los saldos de caja \parencite{Pigou1917Value} y lleva el excedente al bienestar (\ref{sec:interes}).
\item Su \emph{Theory of Unemployment} \parencite{Pigou1933Theory}, en la que el paro se cura bajando los salarios, es el blanco de \autor{Keynes}, que lo llama «clásico» \vertema{3.A.4}.
\item Responde a \autor{Keynes} con el \concepto{efecto \autor{Pigou}}: si bajan los precios, sube el valor real de los saldos monetarios y, con él, el consumo, de modo que el pleno empleo vuelve aunque el tipo de interés no pueda bajar más \parencite{Pigou1943Classical} \vertema{3.A.5}.
\end{itemize}

\begin{anotaciones}
\textbf{\autor{Walras} y \autor{Marshall}, frente a frente}. Los 2 se disputan el título de padre de la economía neoclásica; sus enfoques no compiten, se complementan (tabla~\ref{tab:walras-marshall}).
\begin{tablacaja}
\centering
\caption{\autor{Marshall} y \autor{Walras}: 2 maneras de juntar la oferta y la demanda}\label{tab:walras-marshall}
\small
\begin{tabular}{>{\raggedright\arraybackslash}m{2.6cm}>{\raggedright\arraybackslash}m{5.6cm}>{\raggedright\arraybackslash}m{5.6cm}}
\toprule
\rowcolor{tceeazulpalido} & \textbf{\autor{Marshall}} & \textbf{\autor{Walras}}\\
\midrule
Pregunta & Problemas prácticos: qué pasa con el precio de un bien si cambia su demanda & Una pregunta abstracta: ¿hay unos precios que vacíen todos los mercados a la vez?\\
\addlinespace[0.3em]
Representación & Un mercado tras otro, \emph{ceteris paribus} & Todos los mercados a la vez\\
\addlinespace[0.3em]
Instrumento & Gráficos; las matemáticas, en notas & Sistema de ecuaciones\\
\addlinespace[0.3em]
Ajuste & Por cantidades, según el precio de demanda y el de oferta & Por precios, con el tanteo\\
\addlinespace[0.3em]
Oferta & Depende del plazo: de mercado, corto, largo y secular & Dada en la primera etapa; la producción entra después\\
\addlinespace[0.3em]
Uso hoy & Regular un sector, un impuesto sobre un bien & Una reforma fiscal general, un acuerdo comercial\\
\bottomrule
\end{tabular}
\fuente{elaboración propia a partir de \textcite[cap.~1]{DeVroey2016History}. Cuál conviene depende de lo que se quiera estudiar.}
\end{tablacaja}
\end{anotaciones}

\begin{itemize}
\item \emph{Recapitulación}: la primera generación explica el valor por la utilidad marginal y resuelve la paradoja del agua y los diamantes; Viena lo lleva a los factores, al tiempo y al dinero; Lausana, a todos los mercados a la vez, y Cambridge une la utilidad y el coste en un mercado. Falta todavía llevar el margen a la producción y a la distribución, estudiar con él el interés, el dinero y el bienestar y, cuando la escuela pierda impulso, rehacerla: son la 2.ª y la 3.ª etapa.
\end{itemize}

% ==========================================================================
\section{Productividad marginal y resurgimiento paretiano}\label{sec:segunda}
\minutos{9}

\begin{itemize}
\item Una vez vistos los fundadores, pasamos a la 2.ª y la 3.ª etapa de \autor{Fonseca}: la segunda generación, que lleva el principio del margen a la producción, a la distribución, al interés, al dinero y al bienestar, y el resurgimiento paretiano, que rehace el edificio en los años 30 y 40.
\item \emph{Conclusión anticipada}:
  \begin{itemize}
  \item en los años 90 el margen llega a la oferta: cada factor cobra su productividad marginal y, con rendimientos constantes, el producto se reparte entero;
  \item sobre esa base se levantan las teorías neoclásicas del interés, del dinero y del bienestar;
  \item tras unas décadas de declive, la utilidad ordinal de \autor{Pareto} y los teoremas del bienestar devuelven la escuela al centro de la disciplina, que ya no abandonará.
  \end{itemize}
\end{itemize}

% --------------------------------------------------------------------------
\subsection{Establecimiento de la teoría de la distribución de la productividad marginal (1890-1899)}\label{sec:productividad}

\begin{itemize}
\item Es la 2.ª etapa de \autor{Fonseca}, cuyas fechas son 1890-1894 \parencite{Fonseca2026Neoclassicals}: \textcite{Clark1891Distribution}, \textcite{Wicksteed1894Essay} y \autor{Wicksell} (1893). La obra mayor de \autor{Clark}, \emph{The Distribution of Wealth}, es de 1899.
\item \emph{Las piezas comunes}:
  \begin{itemize}
  \item una \concepto{función de producción}, $q=F(\vect z)$, la mayor producción que se obtiene con los factores $\vect z=(z_1,\dots,z_n)$, con productividades marginales positivas y decrecientes (la ley de las proporciones variables) \vertema{3.A.11};
  \item competencia perfecta en los bienes y en los factores: la empresa toma como dados el precio $p$ y los de los factores $w_i$, y contrata cada factor hasta que el valor de su producto marginal iguala su precio, $w_i=p\,\mathrm{PMg}_i$: la demanda de cada factor es su curva de valor del producto marginal \vertema{3.A.25};
  \item el capital es un factor más, y su remuneración, el interés, se separa del beneficio del empresario, que en equilibrio es nulo.
  \end{itemize}
\item \emph{La pregunta}: los clásicos no necesitaban preguntarse si el producto alcanzaba para todos, porque una de las rentas era el residuo. Ahora cada una se determina por separado: ¿suman el producto?
\end{itemize}

\paragraph{\autor{John Bates Clark} (1847-1938)}

\subparagraph{a) Contexto y obra}

\begin{itemize}
\item Formado 3 años con \autor{Knies} en Heidelberg \parencite[prefacio]{Clark1899Distribution}, introduce el marginalismo en Estados Unidos y es uno de los fundadores de la American Economic Association. En \textcite{Clark1891Distribution} generaliza la renta diferencial de \autor{Ricardo} a todos los factores.
\end{itemize}

\subparagraph{b) La ley de la productividad final}

\begin{itemize}
\item El salario lo fija la productividad del último trabajador, y el interés, la del último capital: «la productividad específica del trabajo fija los salarios» \parencite[pág.~47; traducción propia]{Clark1899Distribution}. Pagar a cada factor lo que produce la unidad marginal «no supone explotación» \parencite[cap.~XXI; traducción propia]{Clark1899Distribution}. «A cada cual lo suyo»: para \autor{Clark}, la competencia da a cada factor lo que aporta.
\item \emph{Estática y dinámica}: en el estado estático, sin cambios en la técnica ni en los gustos, la competencia elimina el beneficio; el beneficio positivo es un fenómeno dinámico, ligado al cambio.
\end{itemize}

\subparagraph{c) Valoración}

\begin{itemize}
\item Su lectura ética es la que más se le discute: la productividad marginal mide lo que se perdería con una unidad menos, no lo que esa unidad «crea», y depende de la cantidad de los demás factores y del reparto de la propiedad.
\end{itemize}

\begin{anotaciones}
\textbf{Más sobre \autor{Clark}}
\begin{itemize}
\item \emph{Influencias}: la escuela histórica alemana, por su formación en Heidelberg, y el socialismo cristiano de su juventud; por eso, en sus primeros escritos, la sociedad y sus leyes (la propiedad, la competencia) pesan tanto como el individuo.
\item \emph{Competencia y poder de mercado}: sus conclusiones dependen de que la competencia funcione. Con monopolio o monopsonio, los factores pueden no cobrar su producto marginal; por eso defiende la competencia y ve en ella la condición del progreso.
\item \emph{Influencia}: en \autor{Fisher}, en la economía estadounidense de la generación siguiente y en la organización industrial. Su hijo, \autor{John Maurice Clark}, será uno de los institucionalistas norteamericanos.
\end{itemize}
\end{anotaciones}

\paragraph{El agotamiento del producto: \autor{Wicksteed} y \autor{Wicksell}}

\subparagraph{a) Idea}

\begin{itemize}
\item \autor{Clark} afirmaba que pagar a cada factor su producto marginal agota el producto, pero no lo demostró. Lo hace \autor{Philip Wicksteed} \parencite*{Wicksteed1894Essay}, que es también uno de los primeros en escribir la producción como una función general de los factores \vertema{3.A.11}.
\end{itemize}

\subparagraph{b) Modelo}

\begin{itemize}
\item \emph{i) Supuestos}: rendimientos constantes a escala (si se duplican todos los factores, se duplica la producción) y competencia perfecta en los bienes y en los factores.
\item \emph{ii) Desarrollo}: por el \concepto{teorema de \autor{Euler}}, con rendimientos constantes la producción es la suma de los factores multiplicados por sus productividades marginales, $\sum_i\mathrm{PMg}_i\,z_i=F(\vect z)$. Con $w_i=p\,\mathrm{PMg}_i$:
  \funcionclave[eq:euler]{\sum_i w_i\,z_i=p\sum_i\mathrm{PMg}_i\,z_i=p\,F(\vect z)}
  Lo que cobran los factores es exactamente el valor de lo producido (anexo~\ref{anexo:derivaciones}).\footnote{\autor{Alfred Flux} (1894), al reseñar el libro, es el primero en ver que el resultado es el teorema de \autor{Euler} \parencite{HET2026Flux}. Con rendimientos crecientes, pagar a cada factor su producto marginal costaría más de lo que se produce; con decrecientes, quedaría un residuo.}
\item \emph{iii) Implicaciones}: no queda residuo y el beneficio extraordinario es nulo: la distribución deja de ser un problema aparte y pasa a ser un problema de precios.
\end{itemize}

\subparagraph{c) Extensiones: \autor{Wicksell}}

\begin{itemize}
\item \autor{Knut Wicksell} (1851-1926) parte de la teoría del capital de \autor{Böhm-Bawerk} y del equilibrio general de \autor{Walras} \parencite{Wicksell1936Interest}. Propone funciones con tramos de rendimientos crecientes, constantes y decrecientes: con libre entrada, cada empresa produce en el mínimo del coste medio, donde los rendimientos son localmente constantes, y el producto se agota aunque la función no tenga rendimientos constantes en todas partes \parencite[págs.~126-131]{Wicksell1934Lectures} \vertema{3.A.12} (anexo~\ref{anexo:derivaciones}).
\end{itemize}

\subparagraph{d) Valoración}

\begin{itemize}
\item \autor{Pareto} y \autor{Edgeworth} criticaron que la demostración dependiera de los rendimientos constantes, y el propio \autor{Wicksteed} acabó retirándola \parencite{Stigler1941Production}.
\item La crítica de fondo llega después: si el capital es un conjunto de bienes distintos, solo se puede sumar en valor, y su valor depende del tipo de interés que se quiere explicar. Es la \emph{controversia del capital} entre las 2 Cambridge, con el \emph{retorno de técnicas} \vertema{3.A.4}.
\item Otras ideas de \autor{Wicksteed}: el coste es siempre un coste de oportunidad, y las relaciones económicas son \emph{no tuistas}: en un intercambio no se persigue el interés de la contraparte (el \emph{tú}), lo que no impide actuar por fines altruistas \parencite{Wicksteed1910Common}.
\end{itemize}

% --------------------------------------------------------------------------
\subsection{Interés, dinero y bienestar}\label{sec:interes}

\paragraph{\autor{Irving Fisher} (1867-1947)}

\subparagraph{a) Contexto y obra}

\begin{itemize}
\item Profesor en Yale, matemático y divulgador, el economista estadounidense más importante de su tiempo. Sus obras se ordenan por temas:
  \begin{itemize}
  \item la utilidad: \emph{Mathematical Investigations} \parencite{Fisher1892Mathematical};
  \item el interés: \emph{Appreciation and Interest} \parencite{Fisher1896Appreciation}, \emph{The Rate of Interest} \parencite{Fisher1907Rate} y \emph{The Theory of Interest} \parencite{Fisher1930Theory};
  \item el dinero: \emph{The Purchasing Power of Money} \parencite{Fisher1911Purchasing} y \emph{The Making of Index Numbers} \parencite{Fisher1922Making};
  \item las crisis: la deflación por deudas \parencite{Fisher1933}.
  \end{itemize}
\end{itemize}

\subparagraph{b) El interés: impaciencia y oportunidad de invertir}

\begin{itemize}
\item \autor{Fisher} recoge 2 de las causas de \autor{Böhm-Bawerk}, que reúne en la \concepto{impaciencia}, y rechaza la 3.ª, la superioridad técnica de los bienes presentes \parencite[cap.~IV]{Fisher1907Rate}: «más vale un toma que dos te daré» es la impaciencia. El interés sale de la impaciencia por gastar la renta y de la oportunidad de invertirla \parencite{Fisher1930Theory}, libro que dedica a \autor{John Rae} y a \autor{Böhm-Bawerk}:
  \begin{equation}
  u'(C_1)=\beta\,(1+r)\,u'(C_2),\qquad F'(K_2)-\delta=r: \label{eq:fisher-interes}
  \end{equation}
  el hogar consume hasta que su impaciencia en el margen iguala el tipo de interés real, y la empresa invierte hasta que el rendimiento neto de la última unidad de capital lo iguala también \vertema{3.A.33}.
\item \emph{El teorema de separación}: con un mercado de capitales perfecto, la decisión de invertir no depende de los gustos de los dueños: primero se elige la inversión que hace máximo el valor actual y después se reparte el consumo \parencite{Fisher1930Theory,Hirshleifer1958Optimal} \vertema{3.A.29}. Anticipa el teorema de \textcite{ModiglianiMiller1958Cost}.
\end{itemize}

\subparagraph{c) El dinero y los precios}

\begin{itemize}
\item \emph{La teoría cuantitativa}: \textcite{Fisher1911Purchasing} escribe la ecuación de cambio,
  \begin{equation}
  M\,V\equiv P\,T, \label{eq:cuantitativa}
  \end{equation}
  con $M$ el dinero, $V$ su velocidad de circulación, $P$ el nivel de precios y $T$ el volumen de transacciones (aquí $T$ no son los impuestos); \autor{Fisher} añade los depósitos con su propia velocidad. Es una identidad, pero se convierte en teoría si $V$ depende de hábitos de pago e instituciones que cambian despacio y $T$ lo fija el sector real: entonces los precios son proporcionales a la cantidad de dinero. Con $k=1/V$, coincide con el enfoque de los saldos de caja de Cambridge (anexo~\ref{anexo:derivaciones}) \vertema{3.A.35}.
  \begin{itemize}
  \item De ahí salen 2 rasgos del modelo neoclásico: la \concepto{neutralidad del dinero}, que solo mueve las variables nominales, y la \concepto{dicotomía clásica}, que estudia por separado los precios relativos (la teoría del valor) y el nivel de precios (la del dinero). El dinero es un velo sobre la economía real.
  \item \autor{Fisher} no lo veía como una propiedad del corto plazo: durante los periodos de transición, mientras se ajustan los tipos de interés, los cambios del dinero sí mueven la producción \parencite[cap.~IV]{Fisher1911Purchasing}.
  \end{itemize}
\item \emph{La ecuación de \autor{Fisher}}, ya en \textcite{Fisher1896Appreciation}: el tipo nominal, $i$, compensa la inflación esperada, $\pi^e$,
  \begin{equation}
  1+i=(1+r)\,(1+\pi^e)\quad\Rightarrow\quad i\simeq r+\pi^e. \label{eq:fisher}
  \end{equation}
  A largo plazo, el tipo nominal sigue a la inflación uno a uno y el real no cambia: es el \concepto{efecto \autor{Fisher}}; a corto, el nominal se ajusta despacio, y el tipo real se mueve \vertema{3.A.41}.
\item \emph{Los números índice}: \textcite{Fisher1922Making} propone el \emph{índice ideal}, la media geométrica de los de \autor{Laspeyres} y \autor{Paasche} \vertema{3.A.9}.
\item \emph{La deflación por deudas}: tras el crac, explica la Gran Depresión por el sobreendeudamiento y la deflación: al liquidar deudas bajan los precios, y con ellos sube el peso real de lo que queda por pagar; «cuanto más pagan los deudores, más deben» \parencite[pág.~344; traducción propia]{Fisher1933}. La idea vuelve en cada crisis financiera \vertema{3.A.42}.
\end{itemize}

\paragraph{\autor{Knut Wicksell}: el tipo natural y el proceso acumulativo}

\begin{itemize}
\item \autor{Wicksell} se pregunta cómo llega el dinero a los precios en una economía con bancos, y responde que a través del tipo de interés \parencite{Wicksell1898Geldzins}.
  \begin{itemize}
  \item El \concepto{tipo natural}, $r^n$, es el que es «neutral respecto de los precios de las mercancías y no tiende a subirlos ni a bajarlos», el que fijarían la oferta y la demanda de capital real si no hubiera dinero \parencite[pág.~102; traducción propia]{Wicksell1936Interest}: el que iguala el ahorro y la inversión \vertema{3.A.33}.
  \item Si el tipo de los bancos, $i$, queda por debajo de $r^n$, la inversión supera al ahorro, el crédito crece y los precios suben, y siguen subiendo mientras dure la diferencia: es el \concepto{proceso acumulativo}. Si queda por encima, bajan.
  \item Como el tipo natural no se observa, propone una regla: subir el tipo si los precios suben y bajarlo si bajan, un antecedente de las reglas de tipos de interés \vertema{3.A.36} y del modelo de la nueva economía keynesiana \vertema{3.A.7}.
  \item De su tipo natural descienden los demás conceptos «naturales» de la macroeconomía: la tasa natural de paro \vertema{3.A.28}, el producto potencial y la brecha de producción. \textcite{Woodford2003Interest} titula su manual de política monetaria con el nombre del libro de \autor{Wicksell}, \emph{Interest and Prices}.
  \end{itemize}
\item De él salen la \emph{escuela de Estocolmo} (\autor{Myrdal}, \autor{Lindahl}, \autor{Lundberg} y \autor{Ohlin}) y, en la hacienda pública, la regla de la unanimidad: cada gasto se aprueba junto con el impuesto que lo financia \parencite{Wicksell1896Finanztheoretische} \vertema{3.A.24}, que \textcite{Lindahl1919Gerechtigkeit} convertirá en sus precios \vertema{3.A.23}.
\end{itemize}

\paragraph{La economía del bienestar: \autor{Edgeworth}, \autor{Pareto} y \autor{Pigou}}

\begin{itemize}
\item \emph{\autor{Francis Ysidro Edgeworth}} (1845-1926), profesor en Oxford y primer director del \emph{Economic Journal}, lleva el cálculo más lejos que nadie en su tiempo \parencite{Edgeworth1881Mathematical}:
  \begin{itemize}
  \item con una utilidad que no es la suma de las de cada bien, dibuja las primeras \concepto{curvas de indiferencia} \vertema{3.A.8};
  \item con la \concepto{curva de contrato}, los repartos en los que ninguno puede mejorar sin que empeore el otro, muestra que en el intercambio bilateral de \autor{Jevons} el precio está indeterminado, y conjetura que la indeterminación desaparece al aumentar el número de agentes: es la idea del \emph{núcleo} \parencite{DebreuScarf1963Limit} \vertema{3.A.21}.\footnote{La caja que hoy lleva su nombre no la dibujó \autor{Edgeworth}, sino \textcite{Pareto1906Manuale}.}
  \item frente al tanteo de \autor{Walras}, imagina el ajuste como una sucesión de contratos provisionales que se rehacen mientras alguien pueda mejorar (el \emph{recontrato}), sin subastador: un antecedente de la teoría de juegos \vertema{3.A.14}.
  \end{itemize}
\item \emph{El óptimo de \autor{Pareto}}: una asignación es óptima si no se puede mejorar a nadie sin empeorar a otro. Es un criterio que no exige comparar utilidades entre personas, y \textcite{Pareto1906Manuale} sostiene que la competencia lleva a él: el germen del primer teorema fundamental de la economía del bienestar (1TFEB) \vertema{3.A.22}.
\item \emph{\autor{Pigou}}: en \emph{Wealth and Welfare} \parencite{Pigou1912Wealth}, ampliado como \emph{The Economics of Welfare} \parencite{Pigou1920Economics}:
  \begin{itemize}
  \item el \emph{bienestar económico} es la parte del bienestar que se puede medir con dinero; aumenta si crece el dividendo nacional y si se reparte mejor, porque, con utilidad marginal decreciente, un euro da más a un pobre que a un rico \vertema{3.A.22};
  \item cuando el producto neto marginal privado se separa del social, el mercado no lleva al óptimo: es la \concepto{externalidad}. Con un coste marginal externo, $\mathrm{CMgS}=\mathrm{CMgP}+\mathrm{CMgE}$, el mercado produce de más, y un impuesto por unidad igual al daño marginal en el óptimo, $t^*=\mathrm{CMgE}(q^*)$, el \concepto{impuesto pigouviano}, lo corrige \vertema{3.A.23}.
  \end{itemize}
\item \emph{Las críticas}: \textcite{Robbins1932Essay} les reprocha comparar utilidades entre personas, lo que lleva a la nueva economía del bienestar, que solo usa el criterio de \autor{Pareto} \vertema{3.A.22}, y \textcite{Coase1960Problem}, a \autor{Pigou}, olvidar que el daño es recíproco y que se puede negociar \vertema{3.A.23}.
\end{itemize}

\begin{anotaciones}
\textbf{\autor{Pareto} más allá del óptimo}
\begin{itemize}
\item \emph{La distribución de la renta}: en el \emph{Cours}, \autor{Pareto} observa que el número de personas con una renta superior a un umbral cae de forma muy regular al subir el umbral, y casi igual de un país a otro: es la \concepto{ley de \autor{Pareto}}, la distribución que hoy se usa para las rentas y los patrimonios más altos \parencite{Pareto1896Cours} \vertema{3.A.45}.
\item \emph{La sociología}: decepcionado con la economía pura, dedica sus últimos años a la sociología, donde describe la circulación de las élites, que se suceden unas a otras en el poder \parencite{Pareto1916Trattato}.
\end{itemize}
\end{anotaciones}

\begin{itemize}
\item \emph{Las piezas juntas}: con precios y salarios flexibles, los mercados se vacían: el paro es voluntario o fruto de rigideces y el dinero solo determina el nivel de precios (teoría cuantitativa). Es el modelo que \autor{Keynes} llamará «clásico» y que se reconstruye mercado a mercado como punto de partida de la \emph{Teoría general} \parencite[cap.~1]{Heijdra2017Foundations} \vertema{3.A.4}.
\item \emph{Recapitulación}: la segunda generación cierra el sistema: reparte el producto según la productividad marginal y explica el interés por la impaciencia y la oportunidad de invertir, los precios por el dinero y el bienestar por el óptimo.
\item \emph{Las críticas}:
  \begin{itemize}
  \item desde fuera llegan de los institucionalistas, que ven en su individuo una calculadora de placeres y dolores \parencite{Veblen1898Why}, de \textcite{Sraffa1926Laws} y de \textcite{Keynes1936};
  \item desde dentro, de los austriacos, para los que el equilibrio deja fuera el tiempo y el conocimiento.
  \end{itemize}
\item Con esas críticas y sin una escuela que las uniera, la teoría pierde impulso a principios del siglo~XX. Pasamos a ver cómo se rehace: la 3.ª etapa.
\end{itemize}

% --------------------------------------------------------------------------
\subsection{Resurgimiento paretiano (1934-1947)}\label{sec:paretiano}

\begin{itemize}
\item Es la 3.ª etapa de \textcite{Fonseca2026Neoclassicals}: una teoría que había perdido impulso se rehace con la utilidad ordinal de \autor{Pareto} y la economía del bienestar, y se convierte en el lenguaje común de la disciplina.
\end{itemize}

\paragraph{a) Contexto: el declive}

\begin{itemize}
\item Para \textcite{Fonseca2026Phases}, la revolución marginalista empieza a perder fuerza en los años 90 y en los 20 está casi parada:
  \begin{itemize}
  \item la utilidad marginal no se puede observar, y a muchos les parece una psicología dudosa;
  \item cada escuela sigue su camino: Cambridge con \autor{Marshall}, Viena, Lausana y Estocolmo, y \autor{Fisher} o \autor{Wicksteed} sin escuela;
  \item en Estados Unidos gana terreno el institucionalismo.
  \end{itemize}
\item El programa de \autor{Pareto} es el único que conserva unos seguidores fieles, sobre todo en Italia, y es el que se recupera en los años 30 \parencite{HET2026Paretian}.
\item \textcite{Robbins1932Essay} prepara el terreno: la economía estudia la elección entre fines y medios escasos, y las comparaciones de utilidad entre personas no son científicas.
\end{itemize}

\paragraph{b) La teoría del valor, con utilidad ordinal}

\begin{itemize}
\item \textcite{HicksAllen1934Reconsideration} derivan la demanda de unas preferencias que solo ordenan las cestas, con curvas de indiferencia y la relación marginal de sustitución, sin medir la utilidad. Recuperan la descomposición de \autor{Eugen Slutsky} \parencite*{Slutsky1915Teoria} del efecto de un precio en un efecto sustitución y un efecto renta \vertema{3.A.8}.
\item \autor{Harold Hotelling} \parencite*{Hotelling1932Edgeworth} deriva la oferta y la demanda de factores de la maximización del beneficio, y \autor{Paul Samuelson} \parencite*{Samuelson1938Note} da a la demanda una base observable, la preferencia revelada \vertema{3.A.9}.
\item \textcite{Hicks1939Value} lleva el enfoque al equilibrio general y a su estabilidad; \textcite{Allais1943Recherche} escribe en Francia un tratado paralelo, y los \emph{Foundations} de \textcite{Samuelson1947Foundations} lo unifican todo: cada resultado es la solución de un problema de optimización o la estabilidad de un equilibrio \vertema{3.A.21}.
\end{itemize}

\paragraph{c) La nueva economía del bienestar}

\begin{itemize}
\item \emph{Los teoremas fundamentales}, que \autor{Pareto} y \autor{Barone} habían sugerido \parencite{HET2026Paretian}:
  \begin{itemize}
  \item el primero dice que todo equilibrio competitivo es un óptimo de \autor{Pareto};
  \item el segundo, que cualquier óptimo de \autor{Pareto} se puede alcanzar como equilibrio competitivo, con una redistribución adecuada de las dotaciones iniciales;
  \item \autor{Abba Lerner} (1934) los muestra con gráficos, y \textcite{Lange1942Foundations} y \textcite{Allais1943Recherche} los demuestran \parencite{HET2026Paretian} \vertema{3.A.22}.
  \end{itemize}
\item \emph{Cómo comparar sin comparar utilidades}:
  \begin{itemize}
  \item los criterios de compensación de \textcite{Kaldor1939Welfare} y \textcite{Hicks1939Foundations}: un cambio mejora si los ganadores podrían compensar a los perdedores \vertema{3.A.22};
  \item la \concepto{función de bienestar social} de \textcite{Bergson1938Reformulation}, que \autor{Samuelson} desarrolla en los \emph{Foundations} \vertema{3.A.24}.
  \end{itemize}
\item \emph{Precio igual a coste marginal}: \autor{Lerner} \parencite*{Lerner1934Concept} y \textcite{Hotelling1938General} lo proponen para los ferrocarriles y los servicios públicos, financiando el déficit con impuestos \vertema{3.A.17}; con los mismos resultados, \textcite{Lange1936Socialism} y \autor{Lerner} defienden un socialismo de mercado frente a \autor{Mises} (\ref{sec:menger}).
\end{itemize}

\paragraph{d) Valoración}

\begin{itemize}
\item Desde los \emph{Foundations}, el método neoclásico, agentes que optimizan y mercados en equilibrio, es el de toda la teoría económica. \textcite{ArrowDebreu1954Existence} sustituyen el cálculo por el análisis convexo y demuestran la existencia del equilibrio de \autor{Walras} \vertema{3.A.21}, y \textcite{Arrow1951Social} muestra con su teorema de imposibilidad los límites de la función de bienestar social \vertema{3.A.24}.
\item La misma generación lo une con \autor{Keynes} en la síntesis neoclásica \vertema{3.A.5}.
\item \emph{Recapitulación}: la 3.ª etapa no inventa el instrumental, que ya estaba en \autor{Edgeworth} y \autor{Pareto}, pero lo libra de la utilidad cardinal y le da una forma axiomática y matemática que ya no abandonará.
\end{itemize}

\begin{anotaciones}
\textbf{La herencia neoclásica en el temario}
\begin{itemize}
\item \emph{Demanda y producción}: la dualidad de \autor{Hotelling}, \autor{René Roy} y \autor{Ronald Shephard} \vertema{3.A.9} \vertema{3.A.12}.
\item \emph{Comercio internacional}: el modelo de \autor{Heckscher}, \autor{Ohlin} y \autor{Samuelson} explica el comercio por las dotaciones de factores \vertema{3.B.5}.
\item \emph{Crecimiento}: el modelo de \autor{Robert Solow} (1956) parte de una función de producción neoclásica con rendimientos decrecientes del capital \vertema{3.A.43}.
\item \emph{Macroeconomía}: la nueva macroeconomía clásica y los modelos de equilibrio general dinámicos y estocásticos vuelven a construir la macroeconomía con agentes que optimizan y mercados en equilibrio \vertema{3.A.6} \vertema{3.A.7}.
\end{itemize}
\end{anotaciones}

\begin{table}[H]
\centering
\caption{Los economistas neoclásicos: obra, escuela e idea fuerza}\label{tab:autores}
\footnotesize
\begin{tabular}{>{\raggedright\arraybackslash}m{2.1cm}>{\raggedright\arraybackslash}m{1.9cm}>{\raggedright\arraybackslash}m{4.3cm}>{\raggedright\arraybackslash}m{4.6cm}>{\centering\arraybackslash}m{1.3cm}}
\toprule
\rowcolor{tceeazulpalido}\textbf{Autor} & \textbf{Escuela} & \textbf{Obra (año)} & \textbf{Idea fuerza} & \textbf{Tema}\\
\midrule
\autor{Jevons} & Británico & \emph{The Theory of Political Economy} (1871) & Grado final de utilidad; ley de la indiferencia & \vertemacorto{3.A.8}\\
\autor{Menger} & \textcolor{tceeverde}{\textbf{Viena}} & \emph{Grundsätze} (1871) & Valor subjetivo; imputación; origen del dinero & \vertemacorto{3.A.35}\\
\autor{Walras} & \textcolor{tceeazul}{\textbf{Lausana}} & \emph{Éléments} (1874-1877) & Equilibrio general; tanteo & \vertemacorto{3.A.21}\\
\autor{Edgeworth} & Británico & \emph{Mathematical Psychics} (1881) & Curvas de indiferencia; curva de contrato & \vertemacorto{3.A.21}\\
\autor{Böhm-Bawerk} & \textcolor{tceeverde}{\textbf{Viena}} & \emph{Positive Theorie des Kapitales} (1889) & 3 causas del interés; rodeos de producción & \vertemacorto{3.A.33}\\
\autor{Wieser} & \textcolor{tceeverde}{\textbf{Viena}} & \emph{Der natürliche Werth} (1889) & Imputación; coste de oportunidad & \vertemacorto{3.A.12}\\
\autor{Marshall} & \textcolor{tceemorado}{\textbf{Cambridge}} & \emph{Principles of Economics} (1890) & Equilibrio parcial; periodos; elasticidad; excedente & \vertemacorto{3.A.16}\\
\autor{Wicksteed} & Británico & \emph{Essay on the Co-ordination of the Laws of Distribution} (1894) & Agotamiento del producto (\autor{Euler}) & \vertemacorto{3.A.11}\\
\autor{Wicksell} & Sueco & \emph{Geldzins und Güterpreise} (1898) & Tipo natural; proceso acumulativo & \vertemacorto{3.A.36}\\
\autor{Clark} & EE.~UU. & \emph{The Distribution of Wealth} (1899) & Productividad marginal de todos los factores & \vertemacorto{3.A.25}\\
\autor{Pareto} & \textcolor{tceeazul}{\textbf{Lausana}} & \emph{Manuale} (1906) & Utilidad ordinal; óptimo de \autor{Pareto} & \vertemacorto{3.A.22}\\
\autor{Fisher} & EE.~UU. & \emph{The Purchasing Power of Money} (1911); \emph{The Theory of Interest} (1930) & Teoría cuantitativa; $i\simeq r+\pi^e$; impaciencia; deflación por deudas & \vertemacorto{3.A.33}\\
\autor{Mises}, \autor{Hayek} & \textcolor{tceeverde}{\textbf{Viena}} & \emph{Theorie des Geldes} (1912); \emph{Prices and Production} (1931) & Ciclo austriaco; cálculo económico; conocimiento disperso & \vertemacorto{3.A.42}\\
\autor{Pigou} & \textcolor{tceemorado}{\textbf{Cambridge}} & \emph{The Economics of Welfare} (1920) & Externalidades; impuesto pigouviano & \vertemacorto{3.A.23}\\
\midrule
\autor{Hicks}, \autor{Allen} & Paretiano & «A Reconsideration of the Theory of Value» (1934) & Utilidad ordinal; efectos sustitución y renta & \vertemacorto{3.A.8}\\
\autor{Lange}, \autor{Allais} & Paretiano & «The Foundations of Welfare Economics» (1942); \emph{À la recherche d'une discipline économique} (1943) & Demostración de los teoremas del bienestar & \vertemacorto{3.A.22}\\
\autor{Samuelson} & Paretiano & \emph{Foundations of Economic Analysis} (1947) & Preferencia revelada; optimización y equilibrio como método & \vertemacorto{3.A.9}\\
\bottomrule
\end{tabular}
\fuente{elaboración propia. La columna «Tema» indica dónde se desarrolla hoy la idea. Debajo de la raya, la 3.ª etapa (\ref{sec:paretiano}).}
\end{table}

% ==========================================================================
\section*{Conclusión}
\minutos{1}

\begin{itemize}
\item \textbf{\emph{Recapitulación}}:
  \begin{itemize}
  \item En conclusión, la exposición deja 3 ideas, una por cada etapa de \autor{Fonseca}:
  \begin{itemize}
  \item el valor depende de la utilidad marginal, no del trabajo, y lo fijan a la vez la oferta y la demanda, en un mercado, con \autor{Marshall}, o en todos, con \autor{Walras}: la revolución de 1871-1874 (\ref{sec:revolucion});
  \item con la productividad marginal, el margen llega a la distribución, y sobre él se levantan las teorías del interés, del dinero y del bienestar (\ref{sec:productividad}, \ref{sec:interes});
  \item tras unas décadas de declive, la utilidad ordinal y los teoremas del bienestar rehacen el edificio entre 1934 y 1947 (\ref{sec:paretiano}).
  \end{itemize}
  \end{itemize}
\item \textbf{\emph{Relevancia}}:
  \begin{itemize}
  \item El instrumental de la microeconomía de hoy (el margen, la optimización y el equilibrio) es el suyo.
  \item El tipo natural y el impuesto pigouviano siguen en la política económica.
  \end{itemize}
\item \textbf{\emph{Extensiones}}:
  \begin{itemize}
  \item el equilibrio general de \textcite{ArrowDebreu1954Existence} \vertema{3.A.21} y la competencia imperfecta de 1933, tras \autor{Sraffa} \vertema{3.A.18};
  \item \autor{Keynes} y los «clásicos» \vertema{3.A.4}, y la síntesis neoclásica \vertema{3.A.5};
  \item la racionalidad limitada y la economía conductual, que ponen en duda al individuo que optimiza \vertema{3.A.1}.
  \end{itemize}
\item \textbf{\emph{Opinión}}:
  \begin{itemize}
  \item Como en toda escuela, algunos programas no tuvieron continuidad: las manchas solares de \autor{Jevons} o el periodo medio de producción de \autor{Böhm-Bawerk}.
  \item 2 críticas siguen vivas: la de Cambridge a la productividad marginal del capital agregado \vertema{3.A.4} y la austriaca al equilibrio estático \parencite{Hayek1948Meaning}.
  \item Con todo, el instrumental ha sobrevivido a sus críticos, porque une en una sola teoría el valor y la distribución.
  \end{itemize}
\item \textbf{\emph{Idea final}}:
  \begin{itemize}
  \item ¿Qué determina el valor, lo que cuesta o lo que se valora? Las 2 cosas, como las 2 hojas de unas tijeras, y siempre en el margen.
  \item \autor{Mill} creía cerrada la teoría del valor en 1848; 23 años después, la paradoja del agua y los diamantes la reabrió.
  \item La respuesta neoclásica sigue siendo la de los manuales y la del refrán: nadie sabe lo que vale el agua hasta que el pozo se seca, porque el valor de una cosa no está en lo útil que es en total, sino en lo que importa su última unidad, y eso depende de lo escasa que sea.
  \item Si la escuela clásica da origen a la economía como ciencia, los neoclásicos la consolidan como tal.
  \end{itemize}
\end{itemize}

\clearpage
\appendix

\section{Anexo: derivaciones}\label{anexo:derivaciones}

\subsection{El intercambio de \autor{Jevons}}\label{anexo:jevons}

\begin{itemize}
\item $A$ tiene $\omega^A_1$ unidades del bien 1 y $B$, $\omega^B_2$ del bien 2. $A$ entrega $x_1$ y recibe $x_2$; por la ley de la indiferencia, todas las unidades se cambian a la misma razón, $p_1/p_2=x_2/x_1$.
\item $A$ elige cuánto entregar:
  \[
  \optim{\max}{x_1}{U^A\big(\omega^A_1-x_1,\ \tfrac{p_1}{p_2}\,x_1\big)}{0\leq x_1\leq\omega^A_1}
  \]
  Con solución interior, $-U^A_1+\tfrac{p_1}{p_2}\,U^A_2=0$, es decir, $U^A_1/U^A_2=p_1/p_2$. Lo mismo para $B$, que entrega $x_2$ y recibe $x_1$: $U^B_1/U^B_2=p_1/p_2$.
\item Igualando, $\mathrm{RMS}^A_{12}=x_2/x_1=\mathrm{RMS}^B_{12}$, la ecuación \eqref{eq:intercambio-jevons}. Con la utilidad aditiva de \eqref{eq:jevons}, cada utilidad marginal solo depende de la cantidad del propio bien.\findemo
\end{itemize}

\subsection{El agotamiento del producto: \autor{Wicksteed} y \autor{Wicksell}}\label{anexo:agotamiento}

\begin{itemize}
\item \emph{Teorema de \autor{Euler}}: si $F(\lambda\vect z)=\lambda F(\vect z)$ para todo $\lambda>0$ (aquí $\lambda$ es el factor de escala, no el multiplicador de \eqref{eq:gossen}), al derivar respecto a $\lambda$, $\sum_i\mathrm{PMg}_i(\lambda\vect z)\,z_i=F(\vect z)$, y en $\lambda=1$, $\sum_i\mathrm{PMg}_i(\vect z)\,z_i=F(\vect z)$ \vertema{3.A.11}.
\item \emph{\autor{Wicksteed}}: la empresa competitiva resuelve
  \[
  \optim{\max}{\vect z}{p\,F(\vect z)-\vect w\cdot\vect z}{\vect z\geq\vect 0}
  \]
  con condiciones de primer orden $p\,\mathrm{PMg}_i=w_i$ en la solución interior. Sustituyendo en el teorema de \autor{Euler}, $\sum_iw_iz_i=p\sum_i\mathrm{PMg}_iz_i=p\,F(\vect z)$, la ecuación \eqref{eq:euler}.
\item \emph{\autor{Wicksell}}: sin rendimientos constantes en toda la función, la elasticidad de escala local en $\vect z$ es $\sum_i\mathrm{PMg}_i\,z_i/F(\vect z)$ \vertema{3.A.11}.
  \begin{itemize}
  \item Al minimizar el coste de producir $q$, $w_i=\mu\,\mathrm{PMg}_i$, con el multiplicador $\mu$ igual al coste marginal, $\mu=\mathrm{CMg}$ \vertema{3.A.12}. Entonces $C=\sum_iw_iz_i=\mathrm{CMg}\sum_i\mathrm{PMg}_iz_i$, y la elasticidad de escala local es $\mathrm{CMe}/\mathrm{CMg}$.
  \item En el mínimo del coste medio, $\mathrm{CMe}=\mathrm{CMg}$ y la elasticidad de escala vale 1: los rendimientos son localmente constantes.
  \item Con competencia y libre entrada, $p=\mathrm{CMg}=\mathrm{CMe}$, y $\sum_iw_iz_i=C=\mathrm{CMe}\,q=p\,q$: el producto se agota aunque la función no tenga rendimientos constantes en todas partes.
  \end{itemize}\findemo
\end{itemize}

\subsection{La ecuación de \autor{Fisher} y los saldos de caja de Cambridge}\label{anexo:fisher}

\begin{itemize}
\item \emph{Ecuación de \autor{Fisher}}: quien presta 1 euro hoy recibe $1+i$ euros mañana, que compran $(1+i)/(1+\pi^e)$ bienes; para que el rendimiento real sea $r$, $1+r=(1+i)/(1+\pi^e)$, es decir, $1+i=(1+r)(1+\pi^e)$. Desarrollando, $i=r+\pi^e+r\,\pi^e\simeq r+\pi^e$, porque el producto de 2 tasas pequeñas es despreciable.
\item \emph{Cambridge y \autor{Fisher}}: con la ecuación de cambio en términos de renta, $M\,V=P\,Y$, y la demanda de Cambridge, $M^D=k\,P\,Y$, el equilibrio del mercado de dinero, $M=M^D$, da $M=k\,P\,Y$; comparando, $V=1/k$. La diferencia está en la lectura: en \autor{Fisher}, $V$ la fijan los hábitos de pago; en Cambridge, $k$ es una decisión de cuánta riqueza guardar en dinero \vertema{3.A.35}.\findemo
\end{itemize}

\section{Anexo: notas del tema antiguo}\label{anexo:notas}

\begin{itemize}
\item Estas notas venían al pie en la versión anterior del tema; no se cantan, pero dan contexto para las preguntas del tribunal.
\end{itemize}

\subsection{Antes de \autor{Bernoulli}: el valor esperado}\label{anexo:pascal}

\begin{itemize}
\item En el siglo~XVII, al nacer la teoría de la probabilidad, matemáticos como \autor{Blaise Pascal} y \autor{Pierre de Fermat} estudiaron cuándo un juego de azar resulta atractivo calculando su valor esperado.
\item La paradoja de San Petersburgo muestra el límite de ese criterio: un juego con valor esperado infinito por el que nadie pagaría mucho. \autor{Bernoulli} la resuelve sustituyendo el valor esperado por la utilidad esperada, con una utilidad marginal del dinero decreciente (\ref{sec:precursores}) \vertema{3.A.10}.
\end{itemize}

\subsection{La escasez de los clásicos y la de los neoclásicos}\label{anexo:escasez}

\begin{itemize}
\item La respuesta neoclásica a la paradoja del agua y los diamantes es que los diamantes valen más porque son más escasos, no porque cueste más producirlos \parencite{Fonseca2026Phases}.
\item Podría parecer lo mismo: si los diamantes cuestan mucho, habrá pocos. \autor{Smith} lo sugiere al explicar el valor de los metales preciosos por su escasez, y su escasez por el trabajo y el gasto que exige extraerlos \parencite[libro~I, cap.~XI]{Smith1776Wealth}.
\item La escasez neoclásica es otra cosa: no basta con que algo sea raro, tiene que ser escaso respecto a lo que se desea. Los huevos podridos de \textcite[cap.~III]{Robbins1932Essay} son raros, pero nadie los quiere; si el deseo de diamantes es muy grande, incluso muchos diamantes son pocos.
\item La teoría neoclásica del valor une, por tanto, 2 ingredientes con una larga historia anterior a 1871:
  \begin{itemize}
  \item el valor relativo de las cosas nace de su escasez relativa;
  \item el deseo subjetivo forma parte de esa escasez.
  \end{itemize}
\item Antes, unos autores relacionaban la rareza con el valor sin preguntarse si era subjetiva, y otros veían la utilidad y la demanda en los precios sin conectarlas con la escasez \parencite{Fonseca2026Phases}.
\end{itemize}

\subsection{\autor{Jevons} y las comparaciones entre personas}\label{anexo:comparaciones}

\begin{itemize}
\item \autor{Jevons} niega que se puedan comparar las utilidades de 2 personas (\ref{sec:ideas-comunes}).
\item En la práctica, sin embargo, a veces razona como si se pudiera: cuando da por supuesto que una renta adicional aporta menos a quien ya tiene mucho que a quien tiene poco, que es el argumento que \autor{Pigou} usará para redistribuir (\ref{sec:interes}).
\end{itemize}

\subsection{Rendimientos a escala y agotamiento del producto}\label{anexo:rendimientos}

\begin{itemize}
\item \emph{Rendimientos crecientes} (función homogénea de grado mayor que 1): los costes medios decrecen y el coste marginal queda por debajo del medio.
  \begin{itemize}
  \item Una empresa que pusiera el precio igual al coste marginal tendría pérdidas.
  \item El producto marginal de cada factor supera a su producto medio, y pagar a cada uno su producto marginal costaría más que lo producido.
  \end{itemize}
\item \emph{Rendimientos decrecientes} (grado menor que 1): los costes medios crecen y el coste marginal supera al medio.
  \begin{itemize}
  \item Con el precio igual al coste marginal, la empresa tiene beneficios.
  \item Pagar a cada factor su producto marginal deja un residuo: los pagos son menores que el producto.
  \end{itemize}
\item Solo con rendimientos constantes, o en el mínimo del coste medio como en \autor{Wicksell}, el producto se agota exactamente (anexo~\ref{anexo:agotamiento}).
\end{itemize}

\subsection{El tipo natural no se observa}\label{anexo:tiponatural}

\begin{itemize}
\item Un tipo de interés por encima del natural frena la actividad y baja los precios, y uno por debajo, lo contrario (\ref{sec:interes}).
\item El problema es que el tipo natural no se observa: hay que estimarlo, y las estimaciones cambian mucho con el modelo \parencite{HolstonLaubachWilliams2017Measuring}. Por eso \autor{Wicksell} proponía guiarse por los precios.
\end{itemize}

\section{Anexo: aspectos para el test}\label{anexo:test}

\begin{anotaciones}
Preguntas tipo test del primer ejercicio: \href{https://www.victorgutierrezmarcos.es/oposicion/temario/primer-ejercicio/test/simulador.html}{simulador de test}.
\end{anotaciones}

\bibliografia
\end{document}
